% 本文件由 LaTeX 排版 Agent 根据有效 translation.json 自动生成。
% author=Sulpitius Severus (c. 363-c. 420); work_id=3503; title=对话集

\chapter{对话集}

% structure: p0001 source=h1 target=omitted reason=page-title-already-rendered-by-parent

对话一\newline{}对话二\newline{}对话三\par

\SourceRecord{Translated by Alexander Roberts.；Nicene and Post-Nicene Fathers, Second Series；Vol. 11.；Edited by Philip Schaff and Henry Wace.；Buffalo, NY: Christian Literature Publishing Co.,；1894.；<http://www.newadvent.org/fathers/3503.htm>.}

% structure: p0001 source=h1 target=section reason=page-title

\section{对话录一}

% structure: p0002 source=h2 target=subsection reason=relative-heading-rank; base=h2

\subsection{关于东方隐修士的德行}

% structure: p0003 source=h2 target=subsection reason=relative-heading-rank; base=h2

\subsection{第一章}

当我和一位高卢朋友聚集在一处时——这位高卢人因他对\Person{玛尔定}的怀念（他曾是\Person{玛尔定}的门徒之一）以及他自身的功德，是我非常亲爱的人——我的朋友\Person{波斯图米亚努斯}加入了我们。他刚为我而从东方归来，三年前他离开祖国去了东方。我拥抱了这位至亲的朋友，亲吻了他的双膝和双脚，我们有一两下仿佛惊愕不已；流着彼此喜悦的泪水，我们走了好一会儿。但不久我们坐在地上铺的麻衣上。然后\Person{波斯图米亚努斯}将目光转向我，先开口说道——\par

\Quote{当我身在\Place{埃及}的偏远地区时，我渴望一路走到海边。我在那里遇到一艘商船，准备启航前往\Place{纳尔榜}。当夜，你仿佛在梦中站在我身边，用手拉住我，带我上那艘船。不久，当黎明驱散黑暗，我在休息的地方起身，心中思量梦境时，突然被对你的渴望抓住，以致毫不延迟地上了船。第三十天在\Place{马赛}登陆，我从那里来，第十天到达这里——如此顺遂的航行赐给了我见你的殷切渴望。只有你，为了你我航过了这么多海，穿越了这么多陆地，请将自己交给我，让我拥抱、享受，与所有人隔绝。}\par

\Quote{我实在，}我说，\Quote{当你还在\Place{埃及}时，我一直在心意思念中与你共融，我对你的情感完全占据了我，因为我日夜默想你。那么，你当然不能想象我现在会一刻也不愿喜悦地注视你，就像我倾心听你一样。我将听你说，我将与你交谈，同时谁也不许进入我们的隐退处，这是我这偏远小室提供给我们的。因为，我想，你不会介意我们这位高卢朋友在场，他如你所见，全心为你的到来而欢喜，正如我自己一样。}\par

\Quote{很好，}\Person{波斯图米亚努斯}说，\Quote{那位高卢人当然会留在我们的同伴中；我虽然不太认识他，但正因为他被你深深喜爱，他不可能不也为我所亲爱。尤其是，因为他是\Person{玛尔定}学派的；我也不会如你所愿，吝于与你作连贯的谈话，因为我正是为此目的而来，即使冒着冗长的风险，也要回应我亲爱的\Person{苏尔皮提乌斯}的愿望。}——说着，他深情地用双手拉住我。\par

% structure: p0008 source=h2 target=subsection reason=relative-heading-rank; base=h2

\subsection{第二章}

\Quote{真的，}我说，\Quote{你清楚证明了真诚的爱能成就多少事，因为你为了我，航过了这么多海，穿越了这么多陆地，可以说是从太阳升起的东方到它落下的西方旅行。那么，来吧，因为我们在此处僻静之地独处，没有其他事务，觉得应当听你说话；来吧，我求你，给我们讲述你全部漂泊的历史。请告诉我们，如果你愿意，基督的信仰在东方如何兴盛；圣人们享有什么平安；隐修士的习俗如何；以及基督在他仆人身上以何等征兆和奇迹行事。因为确实，在我们这地区和我们所处的环境中，生命本身已成为我们的厌倦，我们乐意听你说，是否基督徒甚至在沙漠中也被容许生活。}\par

\Person{波斯图米亚努斯}回答这些话，宣布说：\Quote{我将照你愿望去做。但我求你首先告诉我，所有我离开时在此地作司铎的人，是否仍像我在出发前所认识的那样。}\par

于是我感叹说：\Quote{我求你，不要询问这类问题；我想，你或像我一样知道它们，或如果你不知道，最好什么也不要听。但我不能不指出，你所询问的那些人，不仅没有比你认识他们时更好，甚至有一个曾是我很好的朋友，我习惯从他的爱中得到一些来自其他人迫害的安慰，他却对我表现出了不应有的不友善。然而，我不愿说更严厉的话，因为我曾视他为朋友，甚至当他被认为是我的敌人时，我也爱他。我只补充一句：当我在心中默默思量这些事时，这个忧伤之源深深折磨了我，我几乎失去了一个既明智又虔诚的人的友谊。但让我们转离这些充满悲伤的话题，不如听你讲，按照你刚才的许诺。}\par

\Quote{就这样吧，}\Person{波斯图米亚努斯}喊道。他说这话时，我们都保持沉默，他将他坐着的麻衣向我挪近一些，就这样开始说。\par

% structure: p0013 source=h2 target=subsection reason=relative-heading-rank; base=h2

\subsection{第三章}

三年前，\Person{Sulpitius}，那时我离开此地，向您告别，从\Place{Narbonne}启航，第五日我们进入\Place{非洲}的一个港口：承天主的旨意，航行如此顺利。我心中极其渴望前往\Place{Carthage}，拜访与圣人们相关的那些地方，尤其是到殉道者\Person{Cyprian}的墓前朝拜。第五日我们回到港口，启航进入深海。我们的目的地是\Place{Alexandria}；但南风逆吹，我们几乎被吹到\Place{Syrtis}上；然而谨慎的水手们防范着，抛锚停船。\Place{非洲}大陆便展现在我们眼前；我们乘小艇上岸，发现整个地区荒无人烟，我向内地深入，以便更仔细地探查此地。离海岸约三英里，我见到沙地中央有一间小茅屋，其屋顶，用\Person{Sallust}的话说，如同船底龙骨。它贴近地面，铺着结实的厚木板，并非因为惧怕那里会有暴雨（实际上，那里从未听说过有雨），而是因为那个地区风力强劲，若何时只要一丝微风开始吹起，即便天气相当温和，那些陆地上造成的毁坏也比任何海洋更为严重。那里没有植物，种子从不发芽，因为在如此松动的土壤中，干沙随着每一次风的运动而飞扬。但有些海岬从海边后退，阻挡了风，土壤较为坚固，到处长出一些有刺的草，为羊群提供了不错的牧草。居民以奶为生，而那些较有技巧或较富裕的人则食用大麦面包。这是那里唯一茂盛的谷物，因为大麦在那类土壤中生长迅速，通常能免于烈风的摧残。其生长之快，据说播种后第三十日便成熟。但人们没有理由定居在那里，除非所有人都免于纳税。\Place{Cyrene}人的海岸确实最为遥远，毗邻那片位于\Place{埃及}和\Place{非洲}之间的沙漠，\Person{Cato}从前逃离\Person{Caesar}时曾带领军队穿越此地。\par

% structure: p0015 source=h2 target=subsection reason=relative-heading-rank; base=h2

\subsection{第四章}

\Emph{于是我}朝远处看到的那间茅屋走去。在那里，我遇见一位身穿兽皮衣服的老人，正用手转动磨盘。他亲切地向我们致敬并接待了我们。我们向他解释，我们被迫登陆在那海岸，因海面持续汹涌而无法立即继续航行；我们上岸后，如常人本性使然，希望了解该地的特点和居民的习俗。我们补充说，我们是基督徒，主要询问在这荒野中是否有基督徒。于是，他喜极而泣，匍匐在我们脚前，反复亲吻我们，邀请我们祈祷，同时将羊皮铺在地上，让我们坐下。然后他端上一顿真正奢侈的早餐：半块大麦饼。我们是四人，加上他共五人。他还拿来一捆草药，我忘了名称，但像薄荷，叶子丰盛，味道如蜜。我们非常喜欢这植物异常甜美的味道，饥饿感完全满足了。\par

于是我笑了，对我的朋友\Person{Gaul}说：\Quote{什么，\Person{Gaul}，你觉得怎样？你满意一捆草药和半块大麦饼作为五个人的早餐吗？}\par

他极其谦虚，稍显脸红，同时善意地接受我的玩笑，说：\Quote{\Person{Sulpitius}，你的举动符合你的风格，因为你从不放过任何机会就我们贪吃的话题打趣我们。但你强行让我们\Person{Gaul}人像天使那样生活，未免不仁慈；然而，由于我贪吃，我甚至相信天使也习惯吃东西；因为我的胃口如此之大，我甚至不敢独自对付那半块大麦饼。不过，让那个\Place{Cyrene}人对此满足吧，对他而言，总是感到饥饿要么是必要，要么是本性；或者，又让那些我以为因海上颠簸而失去食欲的人满足吧。我们则远离海洋；正如我常向您见证的，用一句话说，我们是\Person{Gaul}人。但与其在这些事上浪费时间，不如让我们的朋友继续讲述那个\Place{Cyrene}人的故事吧。}\par

% structure: p0019 source=h2 target=subsection reason=relative-heading-rank; base=h2

\subsection{第五章}

\Quote{\Emph{确实}}，\Person{波斯图米亚努斯}继续说，我日后必定留意不再提及任何人的刻苦，以免这艰难的榜样完全冒犯我们的\Place{高卢}朋友。我本想也讲述那位\Place{基勒乃}人的晚餐——因为我们与他共住了七天——或其后的一些宴会；但这些事最好略过，免得高卢人以为被嘲弄。然而，次日当一些当地人聚集来拜访我们时，我们发现那位主人竟是一位司铎——他一直极其小心地对我们隐瞒此事。于是我们随他前往教堂，那堂距离约两英里，中间隔着一座山，使我们无法看见。我们发现它是由普通而不值钱的树木建造的，并不比我们主人的茅舍更堂皇，在那茅舍里若不弯腰便无法站立。在询问当地人的习俗时，我们发现他们既不买也不卖任何东西。他们不知欺诈或偷窃为何物。至于金银，世人通常视之为万物中最可欲的，他们既不拥有，也不渴望拥有。因为当我送给那位司铎十枚金币时，他拒绝了，以深沉的智慧宣称，教会并不因金银而受益，反而受害。但我们仍赠送了他一些衣物。\newline{}\par

% structure: p0021 source=h2 target=subsection reason=relative-heading-rank; base=h2

\subsection{第六章}

在他欣然接受了我们的赠礼后，水手们召唤我们返回大海，我们便离开了；经过一次顺利的航行，我们在第七天抵达了\Place{亚历山大里亚}。在那里我们发现主教与隐修士们之间正爆发一场可耻的争斗，其起因或缘由于此：司铎们被得知常在拥挤的教务会议上通过法令，规定任何人不得阅读或拥有\Person{奥力振}的书。他无疑被视为圣经最有能力的论战者。但主教们坚持认为，在他的书中有些内容是错误的；他的支持者因不敢大胆为这些辩护，便转而宣称它们是被异端者插入的。因此他们断言，不应因那些恰受谴责的部分而谴责他著作的其他部分，因为读者的信德容易加以分辨，使他们不追随伪造的内容，却仍持守那些符合天主教信仰的处理之处。他们指出，在现代和近期的著作中有异端的诡诈作祟，这毫不为奇；因为它甚至在有些地方胆敢攻击福音真理。主教们竭尽全力反对这些主张，坚持认为在\Person{奥力振}的著作中完全正确的部分，连同作者本人，甚至他全部的作品，都应被谴责，因为教会所接受的著作已经绰绰有余。他们还指出，应当避免阅读那些对愚者造成的损害多于对智者益处的著作。至于我，因好奇而查阅了这些著作的一些段落，发现了很多令我喜悦的内容，但也有些该受责备的。我认为显然作者本人确实怀有这些不敬的见解，尽管他的辩护人称那些段落是伪造的。我实在惊奇，同一个人竟能如此自相矛盾：在那被认可的部分，自宗徒时代以来无人能与他相比；而在那被合理谴责的部分，也无人显示比他犯了更严重的错误。\newline{}\par

% structure: p0023 source=h2 target=subsection reason=relative-heading-rank; base=h2

\subsection{第七章}

因为在他的书中，主教们摘录了许多内容，宣读起来表明它们是反对公教信仰而写的，但特别引起对他人反感的一段话是我们在他的著作中读到的：主耶稣，正如祂为救赎人类而降生成人，并为人类的救恩而在十字架上受苦，尝了死味，为人类赢得永生；同样，祂也要通过同一受苦历程，甚至使魔鬼也成为救赎的分享者。他持此论的理由是，这事合乎天主的良善与仁慈，因为祂既恢复了堕落而灭亡的人，也必释放那先前堕落的天使。当主教们提出这些和类似的事情时，由于各方派系的狂热，引起了骚动；而司铎的权威无法平息这骚动，便请来城市长官，以一项乖谬的先例来管理教会的纪律；由于他所引起的恐惧，弟兄们被驱散，隐修士们四处逃窜；以致法令公布后，他们不被允许在任何地方长久地得到接纳。这件事极大地影响了我：热罗尼莫，一个真正公教的人，最精通神圣律法，起初被认为是奥力振的追随者，但现在却超过大多数人，彻底谴责他的全部著作。当然，我不愿对任何人仓促下判断；但据说即使是最博学的人在这场争论中也持有不同意见。然而，奥力振的那个意见，正如我认为的，只是一个错误，还是像普遍认为的那样是一个异端，它不仅不能因司铎们大量的谴责而被压制，而且若没有从他们的争论中汲取力量，它也绝不会传播得如此之广。因此，当我来到亚历山大时，我发现那座城市因与此事有关的骚动而动荡不安。那里的主教非常友善地接待了我，其态度比我所期望的更好，甚至试图挽留我。但我根本无意在那里定居，因为那里最近爆发的恶意导致了弟兄们的毁灭。因为，虽然也许他们似乎应该服从主教们，但如此众多的人，都生活在公开承认基督的状态中，不应因此遭受迫害，尤其不应由主教们来迫害。\par

% structure: p0025 source=h2 target=subsection reason=relative-heading-rank; base=h2

\subsection{第八章}

\Emph{因此}，我从那地方出发，前往伯利恒城，它距离耶路撒冷六英里，但从亚历山大出发需要十六站。司铎热罗尼莫管理此地的教会；因为它是拥有耶路撒冷的主教的一个堂区。我在前一次旅行中已结识了热罗尼莫，他轻易地使我合理地认为没有人比他更值得我尊重和爱戴。因为，除了他因信德和许多德行而应得的功劳外，他不仅精通拉丁文和希腊文，还精通希伯来文，以致没有人敢在所有知识方面与他相比。如果通过他所写的著作您对他还不太熟悉，我将会感到惊讶，因为他实际上被全世界的人阅读。\par

\Quote{好吧，}高卢人此时说道，\Quote{事实上，他太为我们所熟悉了。因为大约五年前，我读了他的一本书，他在书中极其猛烈地攻击和辱骂了我们整个隐修士群体。为此，我们的比利时朋友常非常生气，因为他说我们习惯于填饱肚子甚至到过饱的地步。但我本人宽恕这位杰出的人；我认为他的评论更像是针对东方隐修士而非西方隐修士。因为贪爱饮食在希腊人那里是贪饕，而在高卢人那里则出自他们的本性。}\par

于是我喊道：\Quote{您用修辞学为您的民族辩护，高卢朋友；但我想问，那本书是否只谴责隐修士的这一恶习？}\par

\Quote{不，当然不是，}他回答；\Quote{作者没有放过任何一件事，他指责、鞭笞并揭露一切：他特别抨击了贪婪，也毫不逊色地抨击了傲慢。他大谈骄傲，也谈了不少关于迷信的事；我乐意承认，在我看来他描绘了许多人的恶习的真实图像。}\par

% structure: p0030 source=h2 target=subsection reason=relative-heading-rank; base=h2

\subsection{第九章}

\Quote{至于发生在贞女与隐修士，甚至圣职人员之间的亲昵行为，他的话是多么真实和勇敢！正因如此，据说他不受某些我不愿点名的人的好感。因为，正如我们的比利时朋友因我们被指责过于贪爱饮食而生气那样，那些人又据说当他们在那本小书中读到——\InnerQuote{贞女轻视她真正的未婚兄弟，却寻找一个陌生人}——时，便表达了愤怒。}\par

对此我喊道：\Quote{您说得太过分了，高卢朋友：当心，也许某个承认这些事情的人会听到您的话，并开始对您和热罗尼莫都不太有好感。因为，既然您是一位博学之人，我将不无理由地用那位喜剧诗人的诗句劝诫您——\InnerQuote{顺从赢得朋友，真理招致仇恨。}那么，波司杜米亚努斯，请让我们继续您那开始得很好的关于东方的叙述吧。}\par

\Quote{罢了，}他说道，正如我开始叙述的那样，我与\Person{热罗尼莫}同住了六个月，他不断与恶人作战，且持续与不敬神之人的致命仇恨斗争。异端者恨他，因为他从不停止攻击他们；神职人员恨他，因为他谴责他们的生活和罪行。但毫无疑问，所有善人都钦佩并爱戴他；因为那些认为他是异端者的人是神志不清的。让我在这点上说出真相，那就是此人的学问是大公的，他的教导是纯正的。他终日埋头阅读，全心扑在书本上，昼夜不歇；他不停地要么阅读，要么写作。事实上，若不是我心志已定，且事先向天主许诺要先探访前述的旷野，我几乎不愿离开如此伟大的人片刻。于是，我将我的所有财产和整个家族——他们违背我的意愿跟随我，使我陷入困境——交付并托付给他，从而某种程度上卸下重担，恢复行动自由，我便返回\Place{亚历山大里亚}，探望了那里的弟兄们后，便从那里出发前往上\Place{底比斯}，即\Place{埃及}最偏远的边界。因为据说那片旷野广袤的荒原中住着大量的隐修士。然而，若要叙述我所目睹的一切，就会过于冗长；我只轻描淡写地提几点。\par

% structure: p0034 source=h2 target=subsection reason=relative-heading-rank; base=h2

\subsection{第十章。}

在离旷野不远、靠近尼罗河的地方，有许多隐修院。那里的隐修士大多百人聚居；他们的最高准则就是服从院长的命令行事，凡事不凭己意，一切取决于院长的意愿和权柄。如果碰巧有人心中形成了崇高的德性理想，想要退隐旷野独修，他们不敢擅自行动，除非得到院长的许可。事实上，在他们的评价中，首要的德行就是生活在对他人的服从之中。对于那些退隐旷野的人，面包或其他食物由该院长的命令提供。恰好在我到达那里的那些日子里，院长曾送面包给一位退隐旷野的人，此人在离隐修院不到六里的地方为自己搭了一个帐篷。这面包是由两个少年送去的，年长的十五岁，年幼的十二岁。这两个少年在回家途中，遇到了一条异常巨大的毒蛇，但他们毫不见惧；那蛇仿佛被某种旋律迷住，爬到他们脚边，将暗绿色的脖子伏在他们面前。年幼的那位少年用手抓住了它，用衣服裹住，带着它继续赶路。然后，他像胜利者一样进入隐修院，在众弟兄面前，在所有人的注视下，他打开衣服，放出了被囚的野兽，不无几分夸耀。然而，当其他旁观者称赞孩子们的信德和德行时，院长却以更深的洞察力，为了防止他们在如此年幼便被骄傲冲昏头脑，对两人都施了惩罚。他先大大责备了他们，因为他们在公开场合揭示了主通过他们所行的事。他宣称那不应归功于他们的信德，而应归于天主的德能；并补充说他们应学习在谦卑中事奉天主，而不是以标记和奇迹夸耀；因为自知软弱比任何虚荣地炫耀能力都要好。\par

% structure: p0036 source=h2 target=subsection reason=relative-heading-rank; base=h2

\subsection{第十一章。}

当我提到的那个隐修士听说此事——他得知孩子们因遇到蛇而遭遇危险，而且尽管战胜了蛇，却挨了一顿狠打——他恳求院长今后不要再送任何面包或食物给他。如今，那位属基督的人自甘冒饿死的危险已经八天过去了；他的四肢因禁食而干枯，但他专注于天堂的心智并未动摇；他的身体因克苦而消损，但他的信德依然坚定。与此同时，院长受到圣神的劝告，要去探访那位门徒。出于虔诚的关怀，他急于想知道那位忠信之人靠什么维持生命，因为他已拒绝任何人力帮助他满足需要。于是，他亲自出发去弄清情况。当那隐修士远远看见老人向他走来，便跑去迎接；他感谢他的探访，并领他到自己的小室。当他们一同进入小室时，他们看见挂在门柱上的一个棕榈枝编的篮子里装满热面包。首先闻到热面包的香味；但一摸，面包像是刚从烤炉里取出来的。同时，他们认出这面包不是埃及常见的形状。两人都充满了惊奇，承认这礼物是从天上来的。一方面，隐修士宣称此事是由于院长的到来；另一方面，院长却将其归于隐修士的信德和德行；但两人都极其喜乐地擘开了那由天而降的面包。当老人回到隐修院，将所发生的事报告给弟兄们时，所有人的心中都充满了如此的热忱，以至于争先恐后地急于退隐旷野和其神圣的独居生活；而他们声称自己太可怜了，竟然在人群中居住了那么久，在那里必须忍受人的交往。\par

% structure: p0038 source=h2 target=subsection reason=relative-heading-rank; base=h2

\subsection{第十二章。}

\Quote{在这座隐修院里，我看见了两位老人，据说他们已在那里住了四十年，事实上从未离开过。我认为不应略过不提他们，因为我亲耳听见院长本人和全体弟兄作证，谈及他们的德行：其中一位，太阳从未见过他用餐；另一位，太阳从未见他发怒。}\par

对此，那位高卢人看着我喊道：\Quote{但愿你的那位朋友——我不想提他的名字——此刻在场；我非常希望他听到这个榜样，因为我们许多人已深受他那苦涩的怒气所害。不过，据我所知，他最近已宽恕了他的仇敌；在这种情况下，倘若他听到了那人的行为，便会因此树立的榜样而更加坚定他的宽恕之路，并会感到不受怒气影响是卓越的德行。我不否认他有正当的愤怒理由；但战斗愈艰难，胜利的冠冕愈光荣。因此，我想，若蒙允准，我可以这样说，有一个人应受称赞，因为当他忘恩负义的被释奴离开他时，他更多的是怜悯而非咒骂那逃亡者。而且，他甚至对似乎拐走那奴隶的人也不发怒。}\par

于是我说道：\Quote{若非\Person{波斯图米亚努斯}给了我们那个克服怒气的榜样，我本会对逃亡者的离去非常愤怒；既然不允许发怒，一切这类令人烦恼的记忆，都当从我们心中抹去。让我们，\Person{波斯图米亚努斯}，听听你要说什么吧。}\par

\Quote{我愿照你所求而行，}他说，\Quote{苏尔皮提乌斯，我看你们都如此渴望听我讲述。但请记住，我并非毫无得到更大酬报的希望而向你们说话；只要不久轮到我时，你不拒绝我的请求，我就乐意履行你的要求。}\par

\Quote{我们确实，}我说，\Quote{没有什么可以报答你将要给予我们的恩惠，即使没有更大的回报。然而，请尽管吩咐你所想到的任何事，只要你满足我们的愿望，正如你已经开始做的那样，因为你的话语确实带给我们真正的喜悦。}\par

\Quote{我不会吝惜丝毫，} \Person{波斯图米亚努斯}说，\Quote{不让你们的愿望得到满足；既然你们已认识到一位隐士的德行，我将继续向你们讲述更多关于此类人物的一些事迹。}\par

% structure: p0045 source=h2 target=subsection reason=relative-heading-rank; base=h2

\subsection{第十三章}

那么，当我进入沙漠最靠近的部分，距尼罗河约十二英里时，由一位熟悉地形的弟兄作向导，我们来到一位老隐修士的住所，他住在山脚下。那里有一口井，在此地区非常罕见。那隐修士有一头牛，全部工作就是用轮式机械提水。这是取水的唯一方法，因为据说那井深达一千多英尺。那里还有一个菜园，种满各种蔬菜。这也出乎对沙漠的预料，因为那里一切干枯，被烈日炙烤，连最纤细的植物根茎也不生长。但隐修士与他的牛共同劳动，加上他个人的勤勉，给圣者带来了如此幸福的状况；因为他频繁灌溉，使沙地变得肥沃，以致我们看到菜园里的蔬菜奇妙地繁茂成熟。于是，牛和主人都以此为食；圣者还用这丰盛的出产招待我们用餐。在那里，我看到一件事，你们高卢人或许不会相信——一口锅不用火就煮着准备给我们当晚餐的蔬菜：那地方的太阳威力如此之大，足以胜任任何厨师，甚至能烹调高卢式的佳肴。然后，晚餐后，傍晚来临时，主人邀请我们前往一棵棕榈树，他习惯食用其果实，那树离此约两英里。因为那是沙漠中唯一能找到的树，而且即便这些树也很稀少。我不确定这是出于前代的智慧预见，还是土地自然生长。或许天主预知沙漠将来有一天会被圣徒居住，便为祂的仆人预备了这些。因为定居在这些僻静之处的人大都以这类树的果实为生，因为没有其他植物能在这些地方生长。当我们来到那棵棕榈树前，主人善意引领我们至此，在那里我们遇到一头狮子；一看见它，我和我的向导都开始颤抖；但圣者毫不犹豫地走上前去，我们虽然极为恐惧，还是跟随着他。那野兽仿佛受天主命令，谦恭地退后，站在那里注视我们；而隐修士朋友则从低垂的树枝上摘下一些够得着的果实。当他伸出手，装满椰枣时，那猛兽跑到他跟前，如同驯养的动物一样欢快地接受了，吃完了便离去。我们目睹这一切，仍心有余悸，不能不看透在他身上信德的力量何其大，而在我们身上何其微弱。\par

% structure: p0047 source=h2 target=subsection reason=relative-heading-rank; base=h2

\subsection{第十四章}

我们发现另一位同样非凡的人，住在一间仅能容身的小茅屋里。关于他，我们被告知，一只母狼惯常在他用餐时站在他身旁；而且这野兽绝不会轻易被骗，以至于无法在固定时间他进食时与他在一起。据说，这狼在门口等候，直到他将自己简陋晚餐剩下的面包递给她；她惯常舔他的手，然后，她的职责仿佛完成，向他表示敬意后，便离去了。但恰巧，这位圣人护送一位来访的弟兄回家，耽搁了相当长的时间，直到夜晚才返回。与此同时，那野兽在惯常用餐时间出现了。它进入空置的小屋，发现其恩人不在，便开始好奇地四处搜寻，试图找到居住者。恰巧，一个棕榈叶编的篮子挂在近处，里面有五块面包。野兽取走其中一块，吞吃了，然后犯下这恶行后便离去了。隐修士返回时，发现篮子凌乱，面包数量减少。他意识到家中物品损失，并在门槛附近观察到被偷面包的一些碎屑。考虑到这一切，他几乎毫无疑问地知道窃贼是谁。于是，接下来的日子，那野兽没有像往常一样出现（无疑是由于意识到自己的鲁莽行为，羞于面对它所伤害的人），隐修士因失去与它相处的快乐而深感悲伤。最后，通过他的祈祷，野兽被带回，七天后在惯常用餐时间又出现在他面前。但为了让人人都明白它所感到的羞愧——因所行之事而悔恨——它不敢靠得太近，双眼因深深的自卑而低垂，似乎以某种方式乞求宽恕，这是任何人都能理解的。隐修士怜悯它的窘迫，叫它靠近，然后用慈爱的手抚摸它的头；并给了它两块面包而非平时的一块，使这有罪的生物恢复原状；它得到宽恕后，摆脱了痛苦，重新恢复了往日的习惯。看，我恳求您，即使在这个例子中，也彰显了基督的能力：对他而言，一切非理性都是有智慧的，一切本性凶残之物都是温良的。狼履行职责；狼承认偷窃的罪行；狼因羞耻而困惑；被召唤时，它出现；它伸出头让人抚摸；它感知到所赐的宽恕，仿佛因自己的不端行为而感到羞愧——这是您的能力，哦基督——这些是您的奇妙作为。因为的确，您的仆人在您名下所做的一切，都是您的作为；只有在这一点上，我们深感悲痛：当野兽承认您的威严时，有理智的人却未能崇敬您。\par

% structure: p0049 source=h2 target=subsection reason=relative-heading-rank; base=h2

\subsection{第十五章。}

但恐怕有人会认为此事难以置信，我将提及更重大的事。我呼求基督作证，我并未捏造任何事，也不会引用不可靠作者的说法，而是陈述可靠人士向我证实的事实。\par

有许多人住在沙漠中，没有屋顶遮头，人们称他们为独修者。他们以植物根茎为生；不固定定居在任何地方，以免常有访客；夜晚逼到何处，便在何处安身。当时，有两位来自尼特里亚的隐修士，走向一位按这种方式和条件生活的人。他们虽来自完全不同的地区，却因听闻他的德行，且此人原属同一隐修院时曾是他们亲密挚友，故而前去寻访。他们长期多方寻找；终于在第七个月，在靠近孟斐斯的那片遥远荒野中找到了他。据说他已在这孤寂之地居住十二年；他虽回避与所有人交往，却不避见这些朋友；反而依从他们的情谊，与之共处三日。第四日，当他送他们回程走出一段路时，他们望见一头体型异常硕大的母狮迎面而来。那牲畜虽遇见三人，却毫不迟疑地直奔隐修士，俯伏在他脚下，带着一种哀哭悲鸣伏在那里，流露出忧伤与恳求交织的情感。众人见状皆感震惊，尤其那位意识到自己被寻求的人；于是他动身前行，其余人尾随。那野兽时而停步，时而回头，显然示意隐修士跟随它前往。何需多言？我们来到那牲畜的巢穴，不幸的母狮正在此处养育五只已长大的幼狮，它们自母腹出生时双眼闭合，此后便一直盲目。母狮将幼狮一只只从岩穴中衔出，放在隐修士脚前。圣者终于明白这受造物的所求；他呼求天主圣名，伸手触摸幼狮紧闭的眼睛；它们的盲目即刻消失，久违的光明涌入它们睁开的眼睛。于是，那两位弟兄得以见到渴望拜访的隐修士，带着十分珍贵的辛劳回报而归，既蒙允许亲眼目睹如此大能，便目睹了圣人的信德与基督的光荣，日后他们将为此作证。但我还有更奇妙的事要说：那母狮五日后回到这大施恩惠者处，献上一张不寻常兽皮作为礼物。那位圣人常披此皮如披风，并未轻视藉野兽之手接受的礼物；同时，他更以施予者为另一者。\par

% structure: p0052 source=h2 target=subsection reason=relative-heading-rank; base=h2

\subsection{第十六章。}

在那地区还有另一位独修者的著名名号，他住在赛耶内附近的荒野。此人初入荒野时，原想以沙地各处所生的植物根茎为食，这些根茎味道甘美可口；但他不辨药草性质，常采集有毒的。凭味觉区分根茎种类实非易事，因所有根茎同样甘甜，但许多性质不明者含有致命毒素。当他吃下这些后，体内毒素折磨他，五脏六腑剧痛不已，频繁呕吐伴随难忍痛苦，几乎震碎生命要塞；胃完全无力，他对一切赖以存活的食粮皆感恐惧。这样禁食七日后，他濒临死亡，此时一只称为野山羊的野兽来到他面前。他将前一日采集却未敢触碰的一捆植物递给这站在身旁的野兽；但这野兽用嘴拨开有毒的，挑出它知道无害的。那位圣人由此藉其行为学会当吃什么、当弃什么，既免饥饿致死，又免中毒而亡。但若要详述我们亲身所知或从他人听闻的关于旷野居民的一切事实，实为冗长。我特意在这些孤寂之地度过整整一年再加近七个月，然则我更多是他人德行的赞赏者，而非自身尝试展现他们所显出的非凡坚忍。大部分时间我与前述那位有井有牛的老人同住。\par

% structure: p0054 source=h2 target=subsection reason=relative-heading-rank; base=h2

\subsection{第十七章。}

\Emph{我曾拜访}两个\Person{圣安当}的隐修院，现今由他的门徒居住。我也去了那位最受祝福的\Person{保禄}——隐修者之首——曾居住的地方。我看见了\Place{红海}和\Place{西乃山}的山脊，其山顶几乎触及天际，绝非人力所能攀及。据说有一位隐修者住在山中的某处凹地；我长久且多方寻求见他，却未能如愿。他几乎已有五十年完全断绝与人的交往，不穿衣服，身上覆满自生刚毛，却因天主的恩赐，对自己的赤身无知。每当有虔诚者想拜访他，他便会匆忙逃往无路的旷野，回避一切与人相遇。据说在我到访前大约五年，他曾允许一人与他交谈；我相信那人因信德之力而得此恩惠。在他们两人交谈甚多时，隐修者被问及为何如此刻意避开所有人，他回答说，常被同为人者拜访的人，便不能常被天使拜访。因此，一种传言不无理由地传播开来，并被众人接受：那位圣者享有天使的交往。不论如何，我本人离开了\Place{西乃山}，回到\Place{尼罗河}，我看到河的两岸散布着许多隐修院。我注意到，如我所说，隐修士们大多百人聚居；但有时同住一村的修士竟达两三千人。事实上，人们也无理由认为，大群聚居的修士之德行，逊于那些远离人群独居者。这些地方首要的德行，如我所说，是服从。实际上，任何申请加入者，除非先受试探并证明自己，否则不会为隐修院院长接纳；而且他必须明白，此后绝不能拒绝院长的任何命令，无论多么艰巨困难，甚至看似不堪忍受。\par

% structure: p0056 source=h2 target=subsection reason=relative-heading-rank; base=h2

\subsection{第十八章}

我将叙述两个几乎难以置信的服从之奇事，仅此两例，尽管我记忆中尚有许多；但若区区几个事例不足以激发读者效仿同类德行，再多也无益。那么，当某人放下一切俗务，进入一所戒律严格的隐修院，恳求被接纳为会员时，院长便开始向他提出许多考虑——该修会的辛劳是严酷的；他的要求是沉重的，绝非人的耐性能轻易遵从；他应当另寻生活条件较易的隐修院；不应尝试自己无法达成的事。但此人丝毫不为这些威吓所动；反而更加许诺服从，说如果院长命令他步入火中，他绝不拒绝进入。院长接受了这一宣示后，立即着手考验。恰巧附近有一铁器，被大火烧得极热，正在烘焙一些面包；因炉门敞开，火焰迸发，火在炉膛内肆虐。在此关头，院长命令新人进入火中，他毫不犹豫地服从了命令。他即刻进入火焰之中，而火焰因如此大胆的信德展示，立时被征服，在他临近时退却，正如古时著名的希伯来青年所经历的那样。自然被克服，火焰退让；于是那被认为会被烧死的人，却惊讶地发现自己仿佛被清凉的露水所浸润。哦，\Person{基督}，那火焰为何不触及你年轻的兵士？结果，院长不后悔发出了如此严厉的命令，门徒也不后悔服从了所接到的命令。他确实在来的当日，在软弱中受试探，却被发现是完美的；理当有福，理当光荣，在服从中被试炼，藉受苦而得荣耀。\par

% structure: p0058 source=h2 target=subsection reason=relative-heading-rank; base=h2

\subsection{第十九章}

在同一隐修院中，我所要叙述之事据说发生在不久之前。有一个人同样来到同一院长处，请求加入。当服从的首要法则被摆在他面前，他承诺对一切极端之事有恒久的忍耐时，恰巧院长手中拿着一根早已干枯的\OriginalTerm{苏合香}{storax}枝条。院长将枝条插入土中，并给来访者布置了这项工作：他必须继续浇水，直到（违反一切自然结果）那根枯木在沙土中重现生机。于是，这位陌生人在严苛规则的权下，每天用肩膀从几乎两英里外的\Place{尼罗河}取水来浇灌。一年过去，这位工人的辛劳尚未停歇，却毫无成效的希望。然而，服从的恩宠继续在他的劳动中显现。第二年也嘲弄了这位已疲弱的弟兄的徒劳。终于，在第三年即将过去时，工人日夜不停地浇水，那枝条开始显出生命迹象。我亲眼见到从那条小杖长出的树，至今仍枝繁叶茂地立在隐修院庭院中，仿佛为所述之事作证，显示服从受了何等赏报，信德可发挥何等力量。若要完全叙述我所知的与圣人们之德行相关的许多不同奇迹，时日将不够用。\par

% structure: p0060 source=h2 target=subsection reason=relative-heading-rank; base=h2

\subsection{第二十章}

然而，我还要进一步向你们讲述两件非凡的奇事。其中一件将是对可悲虚荣膨胀的显著警示，另一件则是对虚假正义表露的极好防备。\par

有一位圣人，拥有几乎不可思议的能力，能将魔鬼从附魔者身上驱逐，每日行前所未闻的奇迹。因为他不仅在场时，不仅凭话语，就连不在时，也时常藉取自他衣服的线头或他寄出的书信治愈附身者。因此，从世界各地来见他的人络绎不绝，数量惊人。我不提那些身份卑微的人；长官、朝臣和各级法官常常卧于他的门前。最神圣的主教们，也放下他们司铎的尊荣，谦卑地恳求他触摸并祝福他们，他们合情合理地相信，每当触摸他的手和衣服，他们就圣化了，并蒙受了神圣的恩赐。据说他始终完全戒绝各种饮品，至于食物（我悄悄地告诉你，\Person{苏尔皮基乌斯}，免得我们的高卢朋友听见），仅靠六枚无花果干为生。但在此期间，正如圣人因卓越赢得尊荣，同样，虚荣也从对他的尊敬中悄然潜入。当他初次察觉这恶在他身上增长时，他长久而恳切地挣扎以求摆脱，但无论如何努力都无法彻底根除，因为他仍暗中意识到自己受虚荣的影响。众魔处处承认他的大名，而他无法阻止蜂拥而至的人群。与此同时，隐藏的毒药在他心中作祟，他能一声令下将魔鬼从他人身上赶出，却无法清除自己心中隐藏的虚荣念头。于是他热切恳求天主，据说他祈祷说，让魔鬼有权辖制他五个月，使他变得如同他所治愈的那些人一样。我何必多言？那位最有能力的人——他因奇迹与德行闻名于整个东方，人群聚集其门槛，当代最高权贵俯伏其门前——被魔鬼抓住，牢牢锁在锁链中。直到他受尽附魔者常受的一切苦楚，终于在第五个月被释放，不仅脱离了魔鬼，而且（对他更有益、更可取的）摆脱了住在他里面的虚荣。\par

% structure: p0063 source=h2 target=subsection reason=relative-heading-rank; base=h2

\subsection{第二十一章}

\Quote{但当我思量这些事时，便想到我们自身的不幸与软弱。因为我们中间有谁，若被一个卑贱的人恭敬问候，或被一个女人用愚昧而谄媚的言语称赞，不立刻骄傲自满、膨胀虚荣呢？这会导致即使本人没有圣德的意识，却因愚昧者的谄媚或可能的错认而被称为圣人，事实上他便会自视为至圣！再者，若常有礼物送给他，他便声称自己得天主慷慨如此尊荣，以致一切所需之物在他安眠休息时便赐予了他。更进一步，若有任何能力的迹象，即使微小，临到他身上，他便自视不下于天使。即便他并非因行为或卓越超群，而仅仅被立为神职人员，他便立刻加长衣穗，喜受问候，因人们来访而自大，且自己四处游荡。不仅如此，先前惯于步行或至多骑驴的人，如今必得意地骑在口吐白沫的骏马上；从前满足于简陋小屋的人，如今建造高耸的雕花天花板；他建造许多房间；他雕琢门户；他油漆衣柜；他拒绝粗衣，要求柔软的衣物；并向亲爱的寡妇和友好的贞女们发出如下命令：一些人给他织绣花斗篷，另一些人缝制飘拂长袍。但让我们将这些事留给那位有福的\Person{热罗尼莫}以更辛辣的笔触去描述，我们且回到更直接关注的目标上来。}\par

\Quote{嗯，}我们的高卢朋友对此说道，\Quote{我真不知道热罗尼莫还有什么可说的；你用如此简短的篇幅概括了我们所有的行为，以至于我认为，你的这几句话若被善意接受并耐心思考，将大大有助于那些人，使他们将来无需再被热罗尼莫的著作约束。但你不如继续你已开始的话题，照你所说提出一个反对虚假正义的例子；因为实话告诉你，在整个高卢境内，我们没有比这更有害的恶行了。}\par

\Quote{我就这样做，}\Person{波斯图米亚努斯}回答说，\Quote{我也不再让你们处在期待中了。}\par

% structure: p0067 source=h2 target=subsection reason=relative-heading-rank; base=h2

\subsection{第二十二章}

\Quote{\Emph{一位}来自亚细亚的青年，极其富有，出身望族，有妻室和幼子。他曾在埃及担任护民官，在多次对抗布雷米人的战役中，曾深入沙漠部分地区，也见过几位圣人的帐篷，并从真福若望那里听到了救恩之言。他当时没有延迟，立刻显露出对无益的军旅生活及其虚荣荣耀的轻视。他勇敢地进入旷野，短时间内便在每一种德行上成为完美，因此卓著。他能够长久斋戒，谦逊显著，信德坚定，轻易地获得了与古代隐修士同等的美誉。然而，不久之后，一个（来自魔鬼的）念头进入他心中，认为更恰当的是返回本乡，拯救他的独子、家人以及妻子；这无疑比他只满足于从世界中拯救自己、却不顾（实为不虔敬地）忽略亲友的救恩更能蒙天主悦纳。被这种虚假义德的外貌所诱惑，这位隐修士在将近四年之后，舍弃了他的小室和他所献身的终向。但当他到达最近的住有许多弟兄的隐修院时，他们询问他，他就告知了离去的原因和目的。他们所有人，特别是该院的院长，试图阻止他；但他那不幸形成的意图无法从心中根除。于是，怀着不快乐的固执，他离去，并在众人的悲痛中离开了弟兄们。但他刚从他们视线中消失，就被魔鬼附身，口吐带血的涎沫，开始用牙齿撕咬自己的身体。然后，他被弟兄们抬回同一隐修院，但当这污鬼无法被限制在院墙内时，出于可怕的必要，他被戴上铁链，手脚都被捆绑——这对一个逃亡者来说并非不应得的惩罚，因为现在锁链约束了他，而信德之前未能约束。最终，两年后，因圣人们的祈祷，他从污鬼中获得释放，立刻返回了他曾离开的旷野。这样，他自己得到了纠正，也成了他人的鉴戒，使虚假义德的阴影不致欺骗任何人，也使反复无常的性情不致诱使人以无益的变心，放弃他已踏上的道路。现在，你得知这些关于主在祂仆人身上所施行的各种作为，已足够了；其目的在于或激励他人行类似的行为，或阻止他人行某些事。但因我已充分满足了你的耳朵——实际上已比应当的更冗长——现在（说到这里，他转向我）请你履行你欠我的报酬，让我们按你惯常的方式，听你谈论你的朋友玛尔定；因为我对这事的渴望早已被强烈激起了。}\par

% structure: p0069 source=h2 target=subsection reason=relative-heading-rank; base=h2

\subsection{第二十三章}

\Quote{什么，}我回答说，\Quote{关于我的朋友玛尔定，在你所知我出版的那本关于他的生活和德行的书中，不是已经够多了吗？}\par

\Quote{我承认，}波斯图米亚努斯说，\Quote{你那本书从未远离我的右手。因为如果你认得出，看这里——（说着，他展示了藏在衣服里的那本书）——它就在这里。这本书，}他补充道，\Quote{是我行路和海上的同伴；它是我所有漂泊中的朋友和安慰者。但我将向你讲述那本书已经传到了哪些地方，以及几乎地球上没有一个角落不如这幸福事迹的主题被人熟知。保利努斯，一个对你最敬重的人，首先将它带到了罗马城；随后，当整座城市贪婪地抓住它时，我见到书商们为此欣喜，因为没有什么给他们带来更大的利润，没有什么能更快售出或卖得更高价钱。同一本书，在我旅行中远远领先于我，当我来到阿非利加时，已经在整个迦太基普遍被人阅读。只有我提到的那位基勒乃的司铎没有它；但他根据我的描述将内容记了下来。我又何必提到亚历山大里亚呢？因为在那里，它几乎比你自己更广为人知。它传遍了埃及、尼特里亚、特拜德和整个孟斐斯地区。我在沙漠中见到一位老人正在读它；我告诉他我是你的密友之后，他以及许多其他弟兄托付我这件任务：如果我再次访问这个国家，发现你身体安康，我就应强制你补充你在书中所说关于那位圣人之德行的那些细节，那些你在写作时曾略过的部分。那么，来吧，因为我不想要你向我重复你已写下的那些已经足够熟知的事情，就让我和许多人的请求，将那些你在写作时略过的其他点充分陈述出来，这些点，如我相信，是为了避免读者产生厌烦感。}\par

% structure: p0072 source=h2 target=subsection reason=relative-heading-rank; base=h2

\subsection{第二十四章}

\Quote{事实上，\Person{波斯图米亚努斯}，我一直在专心聆听你谈论圣人的卓越，而与此同时，我内心深处却转向了我的朋友\Person{玛尔定}，我清楚地看到，其他不同的人各自所做的一切，都轻易而完全地被那一个人独自完成了。因为，虽然你确实叙述了崇高的事迹，但我真的没有从你口中听到任何（请允许我这样说，绝无冒犯这些圣人之意）玛尔定不及其中任何一点的事。并且，虽然我认为这些人的卓越无人堪与那人的功德相比，但仍应注意这一点：将他与旷野隐修士甚至隐修士同等比较是不公平的。因为他们摆脱了一切阻碍，只有上天和天使为证，显然受到激励去完成令人赞叹的事；而他呢，身处人群与人事交往之中——好争吵的圣职人员、愤怒的主教周围，被几乎每日发生的丑闻困扰，却仍以不可战胜的德行对这些事全然坚定不移，行了甚至那些我们在旷野中所知或曾经知道的隐修士也没有完成的神迹。但是，即使他们行了与他同等的事，又有哪个审判官会如此不公正，以致没有充分理由断定他更强大呢？试想他是一名在不利地形作战的士兵，最终却成了胜利者；同样地，将他们与那些与敌人势均力敌甚至有利交战的士兵相比。结果如何？虽然胜利对所有人来说都一样，但荣耀肯定不能平等。而且，即使你叙述了神奇的事，你还没有说任何一个人曾使死人复活。毫无疑问，在这一点上必须承认，无人可与\Person{玛尔定}相比。}\par

% structure: p0074 source=h2 target=subsection reason=relative-heading-rank; base=h2

\subsection{第二十五章。}

\Quote{因为，如果那埃及人（如你所言）不被火焰所伤是值得赞叹的，\Person{玛尔定}也屡次显示了他对火的能力。如果你提醒我们，隐修士曾制服并屈服了野兽的凶残，那么玛尔定也惯于约束野兽的狂怒和蛇的毒液。但是，如果你提出那位凭言语的权威，甚至借衣裳的线赶逐附有不洁之魔的人来比较，也有许多证据表明玛尔定在这方面并不逊色。不仅如此，即使你求助于那位以自己毛发为衣、被认为蒙天使探访的人，玛尔定却常与天使交谈。此外，他拥有如此不可战胜的精神以对抗虚荣和自夸，以致无人比他更坚决地鄙视这些恶习，尽管他常常在远处治愈那些附有不洁之魔的人，不仅命令朝臣或总督，甚至命令国王。这在他其他德行中确实是很小的事，但我希望你们相信，没有人比他更积极地对抗虚荣以及虚荣的起因和机会。我还要提到一件小事，但不应忽略，因为它关系到一位拥有最高权力的人，却表现出如此虔诚的愿望来尊敬真福玛尔定。我记得，总督\Person{文森提乌斯}，一位杰出的人，在高卢各种德行上最卓越的人，当他有机会到图尔附近时，多次恳求玛尔定允许他留在隐修院里。他提出这个请求时，举了当时常接待执政官和总督的\Person{圣盎博罗削}主教的例子。但玛尔定以更深的判断拒绝了这样做，以免虚荣和骄傲趁机侵入。因此，你们必须承认，在玛尔定身上存在着你所提及的所有那些人的德行，但在他们所有人身上却找不到使玛尔定出类拔萃的德行。}\par

% structure: p0076 source=h2 target=subsection reason=relative-heading-rank; base=h2

\subsection{第二十六章。}

\Quote{你为什么这样}，波斯塔米亚努斯这时喊道，\Quote{对我说话？好像我并不与你持相同见解，而且从未有过同样心意似的。我确实要在我有生之年，只要我还保有理智，永远歌颂埃及的隐修士；我要赞美独修隐修士；我要赞赏旷野隐修士；但我要把玛尔定放在一个独特的位置上：我不敢拿任何一位隐修士与他相比，更不用说主教了。埃及承认这一点：叙利亚和埃塞俄比亚也发现了这一点：印度听说了这一点；帕提亚和波斯知道了这一点；连亚美尼亚也不无知；遥远的博斯普鲁斯也知道；简而言之，那些到访极乐群岛或北冰洋的人都知晓。我们这片土地因此而更加不幸，它竟不配认识这样一位伟人，尽管他就在近旁。不过，我不把一般民众包括在此指责中：只有神职人员，只有司铎们不知道他；他们心怀恶意不愿认识他，并非没有理由，因为一旦他们认识他的德行，就必须承认自己的恶行。我战栗着说出我最近听到的话：一个可怜的人（我不认识他）说你在你那本书里说了许多谎话。这不是人的声音，而是魔鬼的声音；这样受伤害的不是玛尔定，而是信德本身从福音中被夺走了。因为主亲口为玛尔定所行的那类事作证，说它们应由所有信徒来做，谁若不相信玛尔定行了这些事，就是根本不相信基督说了这些话。但那些可怜、堕落、昏睡的人感到羞耻，因为他们自己不能做的事，他却做到了，他们宁愿否认他的德行，也不愿承认自己的无能。不过，让我们赶着处理其他事，把这些人的记忆都抛开吧：倒不如你，正如我长久以来所盼望的，继续讲述玛尔定尚未被述说的功绩。}\par

\Quote{好吧}，我说，\Quote{我认为你的请求更合适向我们的高卢朋友提出，因为他比我更了解玛尔定的更多事迹——门徒不能不知道师父的事——而且他确实不仅欠玛尔定一份情谊，也欠我们两人，因为我已经出版了我的书，而你则向我们讲述了东方弟兄们的事迹。那么，让我们的高卢朋友开始他应作的详细叙述吧：因为，如我所说，他既欠我们一次说话的回礼，而且我相信，他会为他的朋友玛尔定做这一点——他一定不会不情愿地讲述他的事迹。}\par

% structure: p0079 source=h2 target=subsection reason=relative-heading-rank; base=h2

\subsection{第二十七章。}

\Quote{好吧}，高卢人说，\Quote{我这一方，尽管我不胜任这样重大的任务，却因波斯塔米亚努斯上面所引用的服从范例而感到约束，不能拒绝你加给我的这项职责。但当我想到，我，一个高卢人，要在阿基坦本地人面前说话，我担心我有些粗鲁的言谈会冒犯你们过于挑剔的耳朵。不过，你们将听我说话，如同一个愚拙之人，说话毫不矫揉造作或装腔作势。因为如果你们已承认我是玛尔定的门徒，那么也请允许我，在他的榜样庇护下，鄙弃言语的虚饰和辞藻的华丽。}\par

\Quote{当然}，波斯塔米亚努斯回答，\Quote{你可以说凯尔特话，或高卢话，只要你愿意，只要你说到玛尔定。但在我看来，即使你是哑巴，话语也不会缺少，你可用雄辩的嘴唇说到玛尔定，就如匝加利亚的舌头在提到若翰时松开了一样。但实际上，你是个演说家，你狡猾地像演说家一样，先请求我们原谅你的不熟练，因为你的确在口才上卓越。但一个隐修士不应显示这样的机巧，一个高卢人也不应如此狡诈。倒不如开始工作，把你还有要讲的讲出来吧，因为我们在处理其他事情上已经浪费了太多时间；日渐西斜的阴影警告我们，离黑夜降临已没有多长时间了。} 然后，我们大家沉默了一会儿，高卢人这样开始——\Quote{我想我首先必须注意不要重复我们朋友苏尔皮齐乌斯在那本书里所叙述的关于玛尔定德行的那些细节。为此，我将跳过他的早期事迹，那时他是士兵；我也不涉及他作为平信徒和隐修士所做的事。同时，我什么也不讲只是从别人那里听到的，只讲我自己亲眼目睹的事件。}\par

\SourceRecord{Translated by Alexander Roberts.；Nicene and Post-Nicene Fathers, Second Series；Vol. 11.；Edited by Philip Schaff and Henry Wace.；Buffalo, NY: Christian Literature Publishing Co.,；1894.；<http://www.newadvent.org/fathers/35031.htm>.}

% structure: p0001 source=h1 target=section reason=page-title

\section{对话录 第二卷}

\Emph{论\Person{圣玛尔定}的德行}\par

% structure: p0003 source=h2 target=subsection reason=relative-heading-rank; base=h2

\subsection{第一章}

那么，起初我离开学校后，便追随这位圣人。几天后，我们跟着他前往教堂。路上，一个穷人在这寒冬里半裸着遇见他，乞求给他一些衣服。于是，\Person{玛尔定}叫来总执事，命令立刻给那发抖的人穿上衣服。之后，他进入一间密室，独自坐下，这是他惯常的做法——即使是在教堂里，他也为自己保留这份退隐，因为圣职人员获准如此，何况司铎们坐在另一间屋里，要么相互行礼，要么忙于听审事务。但\Person{玛尔定}一直独自退隐，直到按惯例应该向民众施行神圣礼仪的时刻。我不会忽略这一点：他退隐时从不使用椅子；至于在教堂里，从未有人见过他坐在那里，正如我最近见到某人（天主为我作证），看到那景象不免感到羞愧，他坐在高高的宝座上，实际上，那宝座高如王座；但\Person{玛尔定}却坐在一张粗陋的小凳上，就是最卑微的仆人所用的那种，我们高卢乡民称为\Emph{tripets}，而你们有学问的人，或至少那些来自希腊的人，称为\Emph{tripods}。好了，那个偶然遇见的穷人，因为总执事迟迟不给他衣服，就冲进了圣人的密室，抱怨说他没有得到圣职人员的照顾，并痛苦地哀叹他所受的寒冷。没有耽搁：圣人趁别人没注意，悄悄脱下外衣下面的里衣，用这衣服给穷人穿上，叫他上路。然后，过了一会儿，总执事进来照常通知他，说百姓在教堂里等候，他必须前去举行神圣礼仪。\Person{玛尔定}回答说：\Quote{必须先给穷人穿衣}——他是指自己——\Quote{我无法前往教堂，除非穷人得到了衣服。}但总执事不明白实情——\Person{玛尔定}外面披着斗篷，他看不到里面赤身露体——最后开始抱怨说穷人没有出现。\Person{玛尔定}说：\Quote{那已准备好的衣服，拿来给我；那需要穿衣的穷人不会少的。}于是，总执事迫不得已，而且现在很不耐烦，匆匆从最近的店铺买来一件粗糙的衣服，又短又毛糙，只值五枚银币，怒气冲冲地放在\Person{玛尔定}脚前。他喊道：\Quote{看，衣服在这里，但穷人不在。}\Person{玛尔定}毫不动容，叫他到门口稍等，以此获得秘密，然后自己赤身穿上那件衣服，竭力掩盖他所做的事。然而，在愿意隐藏这些事的圣人身上，什么时候能隐藏得住呢？无论他们愿意与否，一切都被显明出来。\par

% structure: p0005 source=h2 target=subsection reason=relative-heading-rank; base=h2

\subsection{第二章}

\Person{玛尔定}穿着这件衣服前去向天主献祭。就在那天——我要叙述一件奇妙的事——当他照常祝福祭台时，我们看见一团火球从他的头上射出，以至于当它高高升起时，火焰产生了一缕异常长的发丝。虽然我们看见这事发生在非常著名的日子，在众多百姓中，只有一位贞女、一位司铎和三位隐修士目睹了这景象；但为什么其他人没有看见，不是我们的判断所能决定的。\par

大约同时，我的叔父\Person{埃凡提乌斯}，一位非常虔诚的基督徒，虽然忙于世务，却开始患重病，生命垂危，他派人去请\Person{玛尔定}。\Person{玛尔定}毫不迟延地赶往他那里；但圣人还未走完两人之间一半路程，病人就体验到了来者的能力，立刻恢复了健康，他亲自迎接我们到来。他再三恳求，挽留\Person{玛尔定}，\Person{玛尔定}本想次日回家；因为在此期间，一条蟒蛇毒击了我叔父家的一个男孩；\Person{埃凡提乌斯}亲自将他扛在肩上，他几乎因毒力而失去生命，将他放在圣人脚下，相信对圣人来说没有不可能的事。此时，蟒蛇已将毒液扩散到男孩的全身：只见他皮肤上所有血管都肿胀起来，内脏像皮囊一样绷紧。\Person{玛尔定}伸手触摸男孩的四肢，将手指放在那小伤口附近，即毒兽灌入毒液之处。那时，真的——我要讲述奇妙的事——我们看见全身各处的毒液迅速聚集到\Person{玛尔定}的手指上；接着，我们看到毒液混着血液从小伤口处流出，就像牧人用手挤山羊或绵羊的乳头时，大量乳汁常会成线流出一样。男孩完全痊愈了。我们被如此惊人的奇迹所震撼；我们承认——因为真理迫使我们如此——天下无人能与\Person{玛尔定}的事迹相比。\par

% structure: p0008 source=h2 target=subsection reason=relative-heading-rank; base=h2

\subsection{第三章}

同样，此后不久，我们与他一同出行，巡视他教区内的各个堂区。他独自先行了一段路，不知有何必要，我们落在后面。此时，一辆满载军士的公家马车正沿大路驶来。当靠近路边的牲口看见身穿粗糙衣服、外披黑色长斗篷的玛尔定时，受惊而向另一边偏转。随后，缰绳缠绕，使那些（如你常见那样）捆绑这些可怜牲畜的长列陷入混乱；好不容易才重新排列，自然耽误了那些急于赶路的人。士兵们因受此阻而恼怒，急忙跳下地来。他们便开始用鞭子和棍棒击打玛尔定；他默不作声，以令人难以置信的忍耐，将背脊交给他们鞭打，这反而更激起了这些恶徒的怒气，因为见他似乎毫无感觉，仿佛蔑视他们。他几乎昏倒在地；我们不久发现他满身血迹，遍体鳞伤。我们立即扶起他，将他放在自己的驴上，一面咒骂那残暴流血之地，一面尽快赶路。此时，士兵们发泄完怒气后回到车上，催促牲口继续前行。但所有牲口都像铜像一般僵立原地不动；尽管主人大声吆喝，鞭声四起，牲口却纹丝不动。这些人于是纷纷用鞭子抽打；事实上，在惩罚骡子时，他们用尽了所有的高卢鞭子。他们又砍来邻近树林的粗大木棍，用力击打牲口；但这双残忍的手仍无效果：牲口依旧像固定在原地的雕像一样站立不动。这些可怜的人不知所措，再也无法掩饰，这些牲畜心中确有某种更高力量在运作，实际上是被神明的干预所阻止。最后，他们醒悟过来，开始询问刚才在同一地点鞭打的是谁；追查之下，从路人得知正是被他们残忍鞭打的玛尔定。这时，他们不幸的原因对所有人都显而易见；他们不再怀疑是因伤害了那人而受阻。于是，他们全速追赶我们；意识到自己的所作所为和应得的惩罚，满面羞愧，泪流满面，用尘土涂抹头和脸，扑倒在玛尔定脚下，求他宽恕，并恳求他允许他们继续前行。他们补充说，仅凭良心就已惩罚足够，深感大地会在此处活活吞没他们，或者更恰当地说，他们失去知觉，理应变成不动的石头，正如他们看到自己的驮兽固定在所站之处一样；但他们恳求并哀求他赦免他们的罪行，允许他们上路。这位真福者在他们赶来之前已经知道他们被拘留，并且已告知我们此事；但他仁慈地宽恕了他们；并恢复他们的牲口，允许他们继续赶路。\par

% structure: p0010 source=h2 target=subsection reason=relative-heading-rank; base=h2

\subsection{第四章}

我常注意到这一点，\Person{苏尔皮基乌斯}，\Person{玛尔定}惯常对你说，当他做主教时，天主赐予他的能力远不如他记得在获得这一职位之前所拥有的那样丰富。如果这是真的，或者更确切地说，既然这是真的，我们可以想象，当他还是隐修士时，他所成就的那些伟业，以及那些没有见证人而独自完成的，是何等伟大；因为我们看到，他做主教时，在众人眼前行了如此伟大的奇迹。无疑，他从前许多事迹是世人知晓、无法隐藏的，但据说那些因他避免自夸而隐藏起来、不让人知道的事，是不可胜数的；因为他超越常人能力，深知自己的卓越，并践踏世俗荣耀，只满足于让上天作他行为的见证。这一点，即便从我们熟知且无法隐瞒的事中也能判断；例如，在他成为主教之前，他使两个死人复活，你的书已相当详尽地记述了这些事，但在他做主教时，他只复活了一个人，令我惊讶的是你竟未提及。我自己就是后一件事的见证人；但很可能你对此事已有充分证明。无论如何，我将按发生的经过向你陈述。不知为何，我们正前往\Place{卡尔努特斯城}。途中，经过一个居民极为众多的村庄时，一大群人出来迎接我们，全是外教人；因为那村庄里无人认识基督徒。然而，由于这位伟人来临的消息，蜂拥而至的人群充满了广阔的原野。玛尔定感到将有某种工作要做；他内心激动，深受感动。他立刻开始向外教人宣讲天主的圣言，这圣言与人的言语截然不同，他常叹息如此众多的人竟然不认识主救主。这时，当难以置信的人群围住我们时，一个妇人，她的儿子刚去世，她伸出双手将尸体呈给这位有福之人，说：\Quote{我们知道你是天主的朋友：请将我的独生子还给我。}其余的群众也加入她，与母亲一同恳求。玛尔定察觉，如他后来告诉我们的，他能显示权能，为拯救那些等待见此奇事的人，便将死者的身体接在自己手中；当众跪下，祈祷完毕后站起，将复活的孩子还给了母亲。于是，全体群众向天欢呼，承认基督为天主，并最终蜂拥至这位有福之人的膝前，真心恳求他使他们成为基督徒。他也没有迟延。正在原野中央，他按手在他们全体之上，使他们成为慕道者；同时，他转向我们说，在那原野上使这些人成为慕道者并非无故，因为那里曾是殉道者被祝圣的地方。\par

% structure: p0012 source=h2 target=subsection reason=relative-heading-rank; base=h2

\subsection{第五章}

\Quote{你胜利了，啊，\Place{高卢}，}\Person{波斯图米亚努斯}说，\Quote{你胜利了，虽然肯定不是胜过我，我反而是\Person{玛尔定}的维护者，并且一直知道并相信关于此人的这一切；但你胜过了所有隐修士与独修者。因为没有人像你的朋友，更确切地说我们的朋友玛尔定那样，掌管各种死亡。而\Person{苏尔皮基乌斯}那里公正地将他比作宗徒和先知，因为他的信德之能，以及由他的德能所成就的作为，都证明他在各方面都肖似他们。但我恳求你继续讲下去，尽管我们已听到无比精彩的事——但继续吧，啊，\Place{高卢}，将你关于玛尔定还剩下的话说出来。因为心灵急切想知道他哪怕最小最平凡的事迹，因为毫无疑问，他最小的行为也超越他人最大的事迹。}\par

\Quote{我愿如此行事，}高卢人答道，但本人并未亲眼目睹我要讲述的事，因为那发生在我成为玛尔定的同伴之前；尽管如此，这一事实却广为人知，借由当时在场的忠信弟兄们所述而传遍天下。嗯，大约就在他初为主教时，他需要拜访朝廷。那时年长的瓦伦提尼安执掌国事。当他知道玛尔定所请求的是他不愿恩准的事时，便下令禁止他进入宫门。除了他本身的冷酷傲慢外，他的妻子阿里安娜也怂恿他如此，使他完全疏远了圣人，不愿给予他应得的尊重。因此，玛尔定一而再地试图与傲慢的君主会面未果，便诉诸他著名的武器——他披上苦衣，撒灰于身，禁食禁水，昼夜不断祈祷。到了第七天，一位天使向他显现，告诉他放心前往皇宫，因为宫门虽对他紧闭，却会自行打开，而皇帝傲慢的心也会软化。因此，玛尔定因显现天使的话而受到鼓舞，信靠他的帮助，前往皇宫。宫门敞开，无人阻拦他进入；于是他进去，最终来到君王面前，无人试图阻止他。但君王远远看见他走近，咬牙切齿地恼怒他竟被允许进入，丝毫没有屈尊起立，直到火焰覆盖了御座，并燃及君王身体的一部分。如此，傲慢的君主被迫离座，极不情愿地起身迎接玛尔定。他甚至多次拥抱那个他原先决意鄙视的人，心态转好，承认他感知到了神能的运作；甚至不等听取玛尔定的请求，便在他开口之前就恩准了他所求的一切。此后，君王常邀请圣人来商议和宴请；最后，在圣人即将离去时，赠与他许多礼物，但这有福之人谨守自己的贫穷，完全拒绝了，如同他在所有类似场合所做的那样。\par

% structure: p0015 source=h2 target=subsection reason=relative-heading-rank; base=h2

\subsection{第六章}

\Quote{既然我们已进入了皇宫，我便将发生在那里的事件串连起来，尽管它们发生在不同时间。而且，我觉得不应忽略一位忠信的王后所显示的敬仰玛尔定的榜样。当时马克西穆斯统治国家，他本是一个终生都值得称赞的人，倘若他能够拒绝士兵在一场非法骚动中强加于他的王冠，或能够避免内战的话。但事实是，庞大的帝国既无法拒受而不陷于危险，也无法不通过战争来维持。他多次召见玛尔定，迎入宫中，以礼相待；他与玛尔定的谈话始终关乎现在与将来之事、忠信者的荣耀和圣人的不朽；与此同时，王后倾听着玛尔定的唇舌，丝毫不逊于福音中提及的那位妇人，以眼泪洗圣人的脚，用头发擦干。玛尔定此前从未被女人触碰过，却无法逃避她的殷勤，更确切地说，是她仆役般的侍奉。她不顾王国的财富、帝国的尊严、王冠或紫袍，只匍匐在地，无法从玛尔定的脚边被拉开。最后她恳求丈夫（说他们二人应劝说玛尔定同意）让所有其他侍从都离开圣人，只由她一人服侍他用餐。连圣人也无法过于固执地拒绝。她亲手准备素淡的餐食；她亲自为他安排座位；摆放桌子；提供洗手水；端上她自己烹制的食物。他进食时，她双眼低垂，像仆人们一样站在远处一动不动，处处表现出侍奉仆人的谦卑和温顺。她亲自为他调和饮料并奉上。餐毕，她收集用过的面包碎屑，以真切的虔诚将这些剩余视作胜过皇家宴席。有福的妇人！通过展示如此伟大的虔诚，她堪于那位从地极来听撒罗满智慧的妇人相比，如果我们仅从历史的字面意义来看。但两位王后的信德应作比较（请允许我这样说，暂且搁置隐含隐秘真理的庄严）：一位如愿听到了智者之言；另一位不仅堪能听到智者之言，而且得以侍奉他。}\par

% structure: p0017 source=h2 target=subsection reason=relative-heading-rank; base=h2

\subsection{第七章}

波斯图米亚努斯答道：\Quote{高卢人啊，听你说话时，我长久以来一直钦佩那位王后的信德；但你的那句话导向何处，即从未有女人曾更亲近地站在玛尔定身旁？因为让我们想想，那位王后不仅站在他身旁，甚至服侍了他。我实在担心，那些自由地与妇女交往的人，可能会某种程度上以那个例子为自己辩护。}\par

那时，高卢人说：\Quote{你们为什么不注意，就如语法学家常教导我们的那样，地点、时间和人物呢？因为只要你们在眼前呈现这样一幅图画：一个人留在皇宫内，受祈祷的恳求，被皇后的信德所逼迫，为时势所必需，尽他所能去释放被囚禁的人，恢复被流放的人，并收回被夺走的财产——你们认为这些事对一位主教来说该有多么重要，以至于为了实现这一切，他要在自己生活的一般严谨方面稍微放松？然而，既然你们认为有些人会滥用所提供的这个榜样，我只想说，如果他们不逊于这个榜样的卓越之处，他们真是有福的。因为让他们考虑事实是这样的：\Person{玛尔定}一生只有一次，且是在他七十岁时，在用餐时受到服侍和伺候——\Emph{不是}由一位自由的寡妇，也不是由一位放荡的贞女，而是由一位皇后，她在丈夫的权下生活，并且她的行为得到丈夫的恳求支持，才被允许如此行。还要注意的是，她并没有与\Person{玛尔定}一同斜靠在餐桌旁享用宴席，甚至不敢参与宴席，而只是服侍他。因此，要学习正确的做法：让一位有夫之妇服侍你，而不是管理你；让她服侍，但不要与你同席；就如我们所读到的\Person{玛尔大}，她服侍了主，却没有被邀请参与宴席；反之，她（\Person{玛利亚}）宁愿单纯地听道，反而比服侍的（\Person{玛尔大}）更受青睐。但在\Person{玛尔定}的事上，所说的那位皇后履行了两部分：她像\Person{玛尔大}一样服侍，又像\Person{玛利亚}一样听道。因此，如果有人想要使用这个榜样，就让他完全遵守它；让原因相同，人物相同，服侍相同，宴席相同——并且这件事一生只发生一次。}\par

% structure: p0020 source=h2 target=subsection reason=relative-heading-rank; base=h2

\subsection{第八章。}

\Quote{\Emph{钦佩地}}，\Person{波斯图米亚努斯}喊道，\Quote{你的这番话确实约束了我们的那些朋友，使他们不敢超越\Person{玛尔定}的榜样；但我坦白告诉你，我相信你的这些话会落进聋子的耳朵。因为如果我们追随\Person{玛尔定}的道路，我们就永远不需要在亲吻事件上为自己辩护，也就能摆脱所有恶意的指责。但正如你常说的，当你被指责过于贪吃时，你会说：\InnerQuote{我们是高卢人}，同样，我们这些居住在这个地区的人，既不会被\Person{玛尔定}的榜样改变，也不会被你的论述改变。但当我们如此久地讨论这些要点时，你，\Person{苏尔皮基乌斯}，为何保持如此固执的沉默？}\par

\Quote{至于我，}我回答，\Quote{我不仅保持沉默，而且长久以来我已决定在这些问题上保持沉默。因为，我曾责备一位打扮华丽、生活放荡的寡妇，还有一位有点不体面地追随某个我所亲爱的年轻人的贞女（虽然我常听到她指责别人如此行为），这使我在所有女性和所有隐修士中激起了如此大的仇恨，以致两方都对我发动了誓死的战争。因此，请安静，我恳求你们，免得连我们现在说的话也增加他们对我的敌意。让我们彻底忘记这些人，让我们回到\Person{玛尔定}身上吧。你，高卢朋友，如你已开始的那样，继续完成你着手的工作吧。}\par

然后他说：我真的已经向你们讲述了这么多事情，我的讲话应该已经满足了你们的愿望；但是，因为我不能拒绝遵从你们的意愿，只要白天继续，我就会继续讲下去。因为，事实上，当我瞥见那为我们床铺准备的稻草时，我心中浮现出一个关于\Person{玛尔定}曾躺卧过的稻草的记忆，一件奇迹与之相关。事情是这样发生的：\Place{克劳迪奥马古斯}是位于比图里吉人和图罗尼人交界处的一个村庄。那里的教堂因圣人们的虔诚而闻名，并因众多圣洁贞女而同样显赫。\Person{玛尔定}经常经过那里，他在教堂的僻静处有一个房间。他离开后，所有的贞女都涌入那个僻静处；她们亲吻那位有福者曾坐过或站过的每一个地方，并把他躺过的稻草分给自己。几天后，其中一个贞女将她收集来为自己祝福的一部分稻草，挂在了一个被谬误之灵所困扰的附魔者的脖子上。没有延迟，比说话还快，魔鬼被驱逐出去，那人痊愈了。\par

% structure: p0024 source=h2 target=subsection reason=relative-heading-rank; base=h2

\subsection{第九章。}

大约同一时期，当\Person{马丁}从\Place{特里尔}返回时，一头被魔鬼骚扰的母牛遇到了他。那头母牛离开自己的牛群，习惯于攻击人，已经用角严重地顶伤了许多人。现在，当它靠近我们时，那些从远处跟随它的人开始大声警告我们提防它。但它愤怒地逼近我们，眼中充满怒火，\Person{马丁}举起手命令那动物停下，它立刻应声而定。见此，\Person{马丁}看见一个魔鬼坐在它背上，便斥责它，喊道：\Quote{离开，你这致命之物；离开这无辜的牲畜，不要再折磨它。}那恶灵服从并离开了。而这小母牛有足够的理智明白自己获得了自由；因为重获平安后，它俯伏在圣人的脚下；\Person{马丁}指示它后，它便走向自己的牛群，比任何绵羊都温顺地加入了队伍。也是在这时候，他身处火焰中却没有被烧灼的感觉；但我觉得没有必要讲述此事，因为\Person{苏尔皮奇乌斯}在那里虽然在他的书中略过了此事，但在寄给当时身为司铎、如今是主教的\Person{尤西比乌斯}的信函中已相当充分地叙述了。我相信，\Person{波斯图米亚努斯}，你或者已经读过这封信，或者如果它对你仍不熟悉，你可以随时从书架上方便地获取。我将只叙述他省略的细节。\newline{}\par

有一次，当他巡视各个堂区时，我们遇到了一群猎人。猎犬正在追逐一只野兔，那弱小动物已经因长时间奔跑而疲惫不堪。当它看到四周广阔的平原无处可逃，且几次差点被捕获时，它试图通过频繁的折返来延缓死亡。这时，蒙福之人以虔诚之心怜悯这生物的危险，命令猎犬放弃追逐，让它安全逃脱。瞬间，它们一听到第一个命令就完全停住了：人们可能会以为它们被绑住，或者更像是被逮捕，从而在脚印中一动不动。这样，通过追逐者像被绑在一起一样止步，野兔安全逃脱了。\newline{}\par

% structure: p0027 source=h2 target=subsection reason=relative-heading-rank; base=h2

\subsection{第十章。\newline{}}

此外，也值得叙述他的一些常用箴言，因为它们都充满了属灵的教诲。他碰巧看到一只刚被剪过毛的羊，便说：\Quote{看，它完成了福音的诫命：它有两件外衣，将其中一件给了那没有的人；因此，你们也应当这样做。}还有，当他看到一个穿着皮衣、寒冷且几乎赤身的猪倌时，他喊道：\Quote{看那被逐出乐园的亚当，他怎样穿着皮衣喂猪；但是，让我们脱去那仍留在这人身上的旧亚当，反而穿上新亚当。}在某个地方，牛吃掉了草场的青草；猪也用鼻子拱起了一些地方；而剩下未受损的部分，则繁茂生长，仿佛被涂上了各种色彩的花朵。他说：\Quote{那被牲畜吃过的部分，虽然并未完全失去草的美，却没有花朵的壮丽，这象征婚姻；那被不洁的动物猪拱过的部分，则呈现了淫乱的丑恶图画；而那未受损伤的剩余部分，则显扬了童贞的荣耀——它青草繁茂，田野的果实丰盈，并以极美的花朵装饰，犹如璀璨的宝石点缀其上。这样的美是蒙福的，堪配天主；因为无物可与童贞相比。因此，那些将婚姻与淫乱相提并论的人严重错误；而那些认为婚姻应与童贞同等的人则极其可悲和愚蠢。明智之人必须保持这种区分：婚姻属于可以被宽恕的事物，童贞指向荣耀，而淫乱必受惩罚，除非其罪过通过补赎被洗净。}\newline{}\par

% structure: p0029 source=h2 target=subsection reason=relative-heading-rank; base=h2

\subsection{第十一章。\newline{}}

有一位士兵在教会内放弃了军旅生活，宣誓成为隐修士，并在旷野远处为自己搭建了一座小屋，似乎立志要过隐修生活。但久而久之，狡猾的仇敌用各种思想骚扰他属灵的天性，致使他改变心意，希望他的妻子（\Person{玛尔定}曾命令她住在女子隐修院内）能与他同住。于是这位勇敢的隐修士去见\Person{玛尔定}，向他透露了心中的想法。但\Person{玛尔定}坚决否认，认为一个女人不能以不合宜的方式再与一个现在已是隐修士、不再是丈夫的男人结合。最后，当这位士兵坚持己见，声称实现他的目标不会带来任何恶果；他只是想拥有妻子陪伴的安慰；而且他不会再回到从前的行径；他还补充说他是基督的士兵，她也在同一服役中宣誓效忠；因此主教应当允许这些圣人——由于信仰而完全忽略性别问题——一同服役时，\Person{玛尔定}（我要向你们重复他的原话）喊道：\Quote{告诉我，你曾上过战场吗？你曾站过战斗行列吗？}他回答说：\Quote{经常；我常常站在战斗行列中，亲临战阵。}对此\Person{玛尔定}说：\Quote{那么，告诉我，在准备作战的队列中，或是已经前进靠近敌军、拔剑作战时，你可曾见过任何女人站在那里或作战？}这时士兵终于感到困惑和脸红，同时感谢他没有允许他听从自己邪恶的计谋，并且没有用严厉的言辞纠正他，而是通过一个与士兵身份相关的真实而合理的类比。\Person{玛尔定}随即转向我们（因为一大群弟兄围在他周围），说：\Quote{不要让女人进入男人的军营，要让士兵的行列保持独立，妇女应住在自己的帐棚里，远离男人的帐棚。因为如果女流之辈混入男兵营中，军队就会变得可笑。让士兵占据行列，让士兵在平原作战，但女人应将自己留在城墙的保护内。她当然也有自己的荣耀：如果丈夫不在时她保持贞洁；而贞洁的首要美德和圆满胜利，就是她不应被人看见。}\par

% structure: p0031 source=h2 target=subsection reason=relative-heading-rank; base=h2

\subsection{第十二章}

我相信，我亲爱的\Person{苏尔皮奇乌斯}，你还记得他多么强调地赞扬那位贞女——她完全从众人的视线中隐退，以至于当\Person{玛尔定}出于职责想要探访她时，她甚至没有让他进入自己的面前。因为当\Person{玛尔定}经过那片小庄园（她数年来贞洁地把自己关在里面）时，他听说了她的信德和卓绝，便绕道前去，作为主教，他想以虔诚的敬意尊崇这样一份杰出恩赐。我们与他同行的人本以为那位贞女会大大欢喜，因为她将得到如此美德的见证，而一位声望如此崇高的司铎竟一反他平常的严厉态度来拜访她。但她并没有放松她对自己施加的极其严格的生活方式——甚至不让自己见\Person{玛尔定}。于是这位有福之人通过另一位妇女收到了她可称赞的道歉，便愉快地离开了那位不让自己被看见或问候的女子的门口。啊！荣耀的贞女，她甚至不容\Person{玛尔定}看她！啊！有福的\Person{玛尔定}，他不认为那拒绝是对自己的侮辱，反而以欣喜的心赞扬她的卓绝，为在那地方仅属罕见的美德榜样而喜乐。当时，夜幕降临，迫使我们停宿在她简陋居所不远处，那位贞女给有福之人送来了一份礼物；\Person{玛尔定}做了一件他从没做过的事（因为他从不接受任何人的礼物或馈赠）：他一点也没有拒绝那位可敬的贞女送来的东西，并宣称她的祝福绝不能被司铎拒绝，因为她实在应被置于许多司铎之前。我恳请贞女们聆听这榜样，好使她们如果决心向恶人关闭门户，甚至对好人也要关上；为了使恶人不得随意接近她们，她们甚至不必害怕将司铎也拒于交往之外。愿普世都留心听这故事：一位贞女不容许\Person{玛尔定}看她。而她拒绝的不是一个普通的司铎，这个女孩拒绝被一个其面容就是观者救恩的人所看见。但除了\Person{玛尔定}，还有哪个司铎不会认为这是对他的侮辱呢？他会对那贞女心中怀有多么恼怒和愤恨？他会把她视为异端者，并决心把她置于诅咒之下。而这样的人一定会多么肯定地更喜欢那些总是设法出现在司铎面前、大摆宴席、与众人同席的贞女，而不是这位有福的灵魂！但我的话语要把我带到何处？必须克制那稍嫌自由的说法，以免冒犯某些人；因为责备的话对不信者无益，而所举的榜样对信者也足够。同时，我希望如此赞扬这位贞女的美德，却仍然认为不应贬低其他贞女的卓绝——她们常常从遥远的地方来见\Person{玛尔定}，因为事实上，为了同样的目的，天使们也常常拜访这位有福之人。\par

% structure: p0033 source=h2 target=subsection reason=relative-heading-rank; base=h2

\subsection{第十三章}

\Quote{但在我现在要叙述的事上，\newline{}我有你，苏尔皮基乌斯（此时他看着我）作为见证人。有一天，我和苏尔皮基乌斯在那里守在玛尔定的门口，已经沉默地坐了几个小时。我们怀着深深的敬畏和恐惧这样做，仿佛我们在天使的帐篷前执行指定的守夜；而同时，他小室的门关着，他不知道我们在那里。期间，我们听到人们交谈的声音，不久我们充满了一种敬畏和惊奇，因为我们不禁察觉到有神圣的事在进行。将近两小时后，玛尔定出来见我们；然后我们的朋友苏尔皮基乌斯（因为没有人比他更习惯与他亲密交谈）开始恳求他，让我们虔敬地询问这问题，告诉我们那种我们都承认感受到的神圣敬畏是什么意思，以及他在小室里与谁交谈。我们补充说，当我们站在门前时，我们无疑听到了微弱的谈话声，但几乎听不懂。然后他经过很久的迟疑（但确实没有什么苏尔皮基乌斯不能从他那里逼问出来，即使违背他的意愿：我要叙述一些有些难以相信的事，但基督为我作证，我没有说谎，除非有人如此不敬，认为玛尔定本人说谎）说：\InnerQuote{我要告诉\Emph{你}，但我恳求你不要对任何人说。\Person{依搦斯}、\Person{德格拉}和\Person{玛利亚}与我一起在那里。}他接着向我们描述每个人的面容和总体外貌。他承认，不仅在那一天，而且经常，他接待她们的来访。他也不否认，宗徒\Person{伯多禄}和\Person{保禄}也相当频繁地被他看见。此外，他习惯于按魔鬼分别来到他面前时的特殊名字斥责它们。他发现墨丘利是一个特别烦恼的原因，而他说朱庇特愚蠢而迟钝。我知道这些事甚至对许多住在同一隐修院的人来说似乎难以置信；我更不能期望所有仅仅听说它们的人都会相信。因为除非玛尔定过着如此无价的生活，并显示了如此卓越，他绝不会被我们视为被赋予了如此伟大的光荣。然而，人性的软弱怀疑玛尔定的工作一点也不奇怪，因为看到今天许多人不相信福音。但我们自己亲身知道和经历，天使经常显现并与玛尔定亲密交谈。与此相关，我要叙述一件小事，确实不重要，但我仍然要陈述。在\Place{尼姆}举行了一次由主教们组成的主教会议，虽然他拒绝参加，但他仍然想知道会上做了什么。碰巧我们的朋友苏尔皮基乌斯当时与他同船，但按照他的习惯，他待在远离其他人的船上一个僻静处。在那里，一位天使向他宣布了主教会议上发生的事。后来，当我们仔细询问会议举行的时间时，我们无疑发现那正是会议的日子，并且主教们在会上所决定的正是天使向玛尔定宣布的事。}\par

% structure: p0035 source=h2 target=subsection reason=relative-heading-rank; base=h2

\subsection{\Emph{第14章}}

\Quote{但当我们问他关于世界末日的事时，他对我们说，尼禄和假基督必须先来；尼禄将在征服十个国王后统治世界的西方部分，他将进行迫害，迫使人们崇拜外邦人的偶像。他还说，假基督另一方面将首先夺取东方的帝国，以耶路撒冷为他的宝座和国都，而城市和圣殿都将被他重建。他补充说，他的迫害将旨在迫使人们否认基督是天主，而他却坚持自己是基督，并命令所有男人按照法律受割损。他进一步说，尼禄将被假基督摧毁，整个世界和所有民族都将被假基督的权力所控制，直到那邪恶者被基督的来临推翻。他还告诉我们，毫无疑问，假基督由恶神受孕，已经出生，并且此时已到了少年时期，而一旦达到适当年龄他就要掌握权力。现在，我们听他口中说出这些话已有八年：你们可以猜测，那些将来所恐惧的事有多么接近发生。}\par

当我们的高卢朋友激动地说着这些，还没有讲完他打算叙述的事时，家中的一个少年进来宣布，司铎雷弗里吉乌斯正站在门口。我们开始犹豫，是继续听高卢朋友讲下去，还是去欢迎那位我们如此深爱、并且来拜访我们的人；这时我们的高卢朋友说：\Quote{即使这位最神圣的司铎没有到达，我们的谈话也不得不中断，因为夜晚的临近本身就在催促我们结束这已经持续了这么久的讲论。但是，既然关于玛尔定卓越之处的所有事情还远没有被提及，今天你们听到的就够了；明天我们继续剩下的。}我们所有的人都同样喜欢这位高卢朋友的承诺，于是站了起来。\par

\SourceRecord{Translated by Alexander Roberts.；Nicene and Post-Nicene Fathers, Second Series；Vol. 11.；Edited by Philip Schaff and Henry Wace.；Buffalo, NY: Christian Literature Publishing Co.,；1894.；<http://www.newadvent.org/fathers/35032.htm>.}

% structure: p0001 source=h1 target=section reason=page-title

\section{对话第三}

\Emph{玛尔定之诸德（续）}\par

% structure: p0003 source=h2 target=subsection reason=relative-heading-rank; base=h2

\subsection{第一章}

\Quote{天已破晓，高卢友，你该起身了。因为如你所见，波斯特米亚努斯早已急不可待，而这位昨天已被允许旁听的司铎，也期待你能兑现诺言，今天就你所推迟讲述的、关于我们亲爱的玛尔定的事迹，继续往下说。他并非不知道将要叙述的一切，然而即使是回顾已知之事，知识也是甘美而愉悦的；实际上，自然如此安排，使人因着众多见证所确证、毫无疑义之事而更安心地享受认知之乐。因为此人自幼追随玛尔定，确实知晓他的一切作为，但他仍乐意重听已知之事。我向您承认，高卢人啊，玛尔定的德行我常有所闻，甚至我已将许多关于他的事写下；但由于对他事迹的崇敬，那些关于他而一再重复的事，对我总是新鲜的。因此，我们为你庆贺，雷夫里热里乌斯也加入我们成为听者，而波斯特米亚努斯愈发热切，因为他仿佛急于将这些事传至东方，如今要从你这里听到由见证者所确认的真理。}\par

当我正说着这些话，高卢人已准备继续叙述时，一群隐修士突然涌入我们中间：司铎埃瓦格里乌斯、阿佩尔、萨巴提乌斯、阿格里科拉；稍后，司铎埃泰里乌斯与执事卡卢皮奥、副执事阿马托尔也进来了；最后，我的一位至友司铎奥雷里乌斯从更远处赶来，气喘吁吁地冲了进来。我问：\Quote{你们为何如此突然且不约而同地从四面八方跑到我们这里，而且是在清晨这么早的时候？}他们答道：\Quote{我们昨天听说，你的高卢朋友用了一整天讲述玛尔定的德行，因夜幕降临而将余下部分推迟到今天；所以我们急忙赶来为他提供众多听众，听他讲述如此有趣的话题。}与此同时，我们得知门外站着一群平信徒，不敢进入，但仍恳求能被允许进入。阿佩尔于是宣称：\Quote{这些人绝不该与我们混在一起，因为他们来听讲，是出于好奇而非虔诚。}我为那些他认为是不能被接纳的人感到难过，但我所能争取到的，且费了很大力气，只是让他们从侍卫官中接纳了欧赫里乌斯，以及一位执政官级人物凯尔苏斯，其余的人则被挡在门外。然后我们将高卢人安排在中间座位；他依照其一贯谦虚保持久久的沉默，终于开始如下所述。\par

% structure: p0006 source=h2 target=subsection reason=relative-heading-rank; base=h2

\subsection{第二章}

虔诚而口才出众的朋友们，你们聚集起来听我讲话；但依我推测，你们带的是宗教的而非博学的耳朵；因为你们将听我作为信德的见证人讲话，而非以演说家的流利口才。我不会重复昨天说过的事：那些没有听到的人可以通过书面记录了解它们。波斯特米亚努斯期待听到新事，意图将他所听到的传至东方，使东方不因与玛尔定比较而自视高过西方。首先，我的心倾向于讲述一件雷夫里热里乌斯刚才在我耳边低声提及的事：事情发生在卡尔努特城。一位一家之主斗胆将他一出生就哑巴的十二岁女儿带到玛尔定面前，恳求这位受祝福的人能通过他虔诚的功绩解开那被缠住的舌头。玛尔定让位于当时恰好在他身边的瓦伦提努斯主教和维克特里基乌斯主教，声称自己不配如此大事，但对他们（那些似乎比自己更圣洁的人）来说没有不可能的事。然而，他们加上虔诚的恳求，用恳求的声音同那位父亲一起，恳求玛尔定完成所期望的事。他便不再耽搁——他在两方面都可钦佩：首先显示谦卑，然后不推诿虔诚的责任——他命站在周围的群众离开；只留下两位主教和那女孩的父亲，他照常伏地祈祷。然后他祝福一点油，同时念出驱魔的公式；用手指捏住女孩的舌头，将祝圣过的油倒入口中。所施展能力的结果并没有令圣洁之人失望。他问她父亲的名字，她立刻回答。父亲喊着，抱住玛尔定的膝盖，喜极而泣；众人惊讶，他承认那是他第一次听到女儿的声音。为使此事不被视为难以置信，让在场的埃瓦格里乌斯为你们提供真实的见证；因为事情就发生在他的眼前。\par

% structure: p0008 source=h2 target=subsection reason=relative-heading-rank; base=h2

\subsection{第三章}

\Quote{以下是近日从司铎\Person{亚尔帕久}的叙述中得知的一件小事，但我想不应略过不提。廷臣\Person{阿维提亚努斯}的妻子曾送一些油给\Person{玛尔定}，请他祝福（这是惯常做法），以便在必要时用以应付各种疾病。油盛在一个玻璃罐中，形状浑圆，中部逐渐鼓出，瓶颈细长；但伸长的瓶颈是空的，因为灌装这类容器时通常要让顶部空出来，以便塞上封口的塞子。司铎作证说，他看见油在\Person{玛尔定}的祝福下增多，多到溢出罐子，从顶部四面八方流下。他还说，当罐子被送回女主人手中时，同样涌溢出来；因为油在搬运的男孩手中不断溢出，流出的油浸湿了他整个衣服。他更说，女主人接到的罐子满到边缘，以至于（正如这位司铎今日告诉我们的）那罐子再没有空隙可以塞入通常用来封闭需特别保存之物的塞子。发生于此人身上的也是一件奇事。} 说到这里，他看了看我。他曾将一个装有\Person{玛尔定}祝福之油的玻璃罐放在一个相当高的窗台上；家中的一个孩子不知道那里有个罐子，用力过猛地拉下了盖罐子的布。于是罐子掉在大理石地板上。当时所有人都充满恐惧，生怕\Person{玛尔定}赐予罐子的天主的祝福丢失了；但罐子安然无恙，仿佛掉在最柔软的羽毛上。如今这一结果不应归因于偶然，而应归因于\Person{玛尔定}的能力，他的祝福不可能失效。\par

还有一件由某人成就的事，因他在场且禁止提及名字，故隐去其名：\Person{撒杜尔尼诺}，他现在与我们同在，也在所述场合在场。一条狗以一种颇为不悦的方式向我们吠叫。那人说：\Quote{我因\Person{玛尔定}之名，命令你安静。}那狗——它的吠叫似乎卡在喉咙里，人们可能以为它的舌头被割掉了——\Emph{是}沉默的。如此看来，\Person{玛尔定}本人行奇迹实在算不得什么：请相信我，其他人也因他的名成就了许多事。\par

% structure: p0011 source=h2 target=subsection reason=relative-heading-rank; base=h2

\subsection{第四章}

你曾知晓前廷臣\Person{阿维提亚努斯}极其野蛮、无比残忍的凶暴。他怒不可遏地进入\Place{图尔}城，成群身戴锁链的人神情忧郁地跟随其后，他命令备好各式刑罚来处决他们；令全城大为惊愕的是，他安排次日执行这悲惨的行径。\Person{玛尔定}得知此事后，大约半夜前，独自前往那畜生的府邸。但在深夜的寂静中，当众人安寝时，大门闩着无法进入，他便躺在那残忍的门槛前。与此同时，\Person{阿维提亚努斯}沉睡正酣，被一位攻击的天使击打，对他说：\Quote{天主的仆人卧在你门槛上，你还在安睡吗？}他听了这话，从床上惊慌起身，呼唤仆人，恐惧地喊道：\Quote{玛尔定在门口：快去打开门闩，免得天主的仆人受到伤害。}但他们，像所有的仆人一样，刚迈出第一道门槛，就嘲笑主人被梦捉弄，断言门口无人。他们这样做只是根据自身的性情推断，认为不可能有人彻夜守望，更不相信一位司铎会在那样恐怖的黑夜躺在别人的门槛上。他们轻易说服了\Person{阿维提亚努斯}相信他们的话。他再次沉入睡眠；但不久又被更猛烈地击打，他喊道玛尔定\Emph{是}站在门口，因此身心都不得安宁。仆人迟疑时，他亲自走到外面的门槛前，果然在那里找到了\Person{玛尔定}，正如他所想。这可怜的人被如此卓越的显现所震撼，喊道：\Quote{先生，你为什么对我这样做？你不必说话：我知道你所愿：我明白你所求：赶快离开，免得天主的愤怒因我伤害你而将我吞噬：我已经受了足够的惩罚。请相信我，我已在我心中坚定决定了现在该如何行事。}于是，圣人离开后，他召来官员，下令释放所有囚犯，随即自己也离开了。这样，\Person{阿维提亚努斯}被赶走，全城欢庆，感到自由。\par

% structure: p0013 source=h2 target=subsection reason=relative-heading-rank; base=h2

\subsection{第五章}

虽然这些都是确实的事情，因为阿维提亚努斯曾向许多人讲述过，然而还有进一步的证实：你们现在在场的司铎雷弗里热里乌斯，最近曾听护民官中一位忠信的人达格里杜斯，在天主威严的呼求下，将这一切叙述给他；达格里杜斯发誓说，这叙述是阿维提亚努斯亲自告诉他的。但我不希望你们惊奇我今天做了昨天没做的事，即我在提及每一件奇事之后，附上了见证人的名字，并提到一些仍然在世的人，若有人倾向于不信，可以去找他们。许多人的不信迫使我这样做；因为有人对昨天讲述的一些事表示犹豫。那么，让这些人接受仍然健在的人为见证人，既然他们怀疑我们的诚意，就让他们更相信那些见证人吧。但实际上，如果他们如此不信，我认为即使提到了见证人，他们也不会相信。然而我感到惊讶的是，任何稍微有点宗教情感的人，竟敢如此邪恶，认为有人会关于玛尔定说谎。愿这远离每一个服从天主的人；因为玛尔定确实不需要谎言来辩护。但是，基督啊，我们将我们全部言论的真实性摆在祢面前，即我们既没有说、也不会说任何不同于我们亲眼所见、或从无可置疑的权威那里所听到的，而且确实非常频繁地从玛尔定本人那里听到的。虽然我们采用了对话的形式，以使文体变化避免厌倦，但我们仍肯定我们是在以孝爱精神陈述一段真实的历史。一些人的不信迫使我非常遗憾地在叙述中插入了这些离题的评论。现在让言论回到我们的集会吧；在那里，既然我看到自己被如此热切地倾听，我觉得有必要承认阿佩尔把不信者挡在门外做得对，因为他确信只应允许有信德的人听闻。\par

% structure: p0015 source=h2 target=subsection reason=relative-heading-rank; base=h2

\subsection{第六章}

相信我，我心里愤怒，而且因恼怒似乎失去了理智：难道基督徒不相信玛尔定的异能，连魔鬼也都承认的吗？\par

那位有福者的隐修院离城有两英里；但是，每当他要到教堂来时，他只要一踏出小室的门槛，就能听见整个教堂里附魔者的吼叫声，有罪者的群体颤抖如审判官来临，以致魔鬼的呻吟向那些先前不知道主教要来的神职们宣告了主教的临近。我看见一个人，在玛尔定接近时被提到空中，双手向上伸展悬在那里，以致双脚完全不能触地。但若玛尔定有时承担驱魔的任务，他不用手触摸任何人，也不用言语斥责任何人，不像神职们通常那样滔滔不绝地念许多咒语；相反，附魔者被带到他面前时，他命令所有其他人离开，门闩上，自己穿着粗毛衣，撒上灰，躺在教堂中央的地上，转面祈祷。那时真能看见那些可怜的人以各种方式受苦——有的像悬在云中，脚朝上，但他们的衣服没有垂到脸上，以免暴露的部位引起羞耻；而在教堂的另一处，能看见他们未经任何审问就受苦，并承认自己的罪行。他们也自愿透露自己的名字：一个承认是朱庇特，另一个承认是墨丘利。最后，能看见所有魔鬼的仆人连同他们的主人一同痛苦，以致我们不得不承认在玛尔定身上应验了经上所记载的：\BibleQuote{圣人将要审判天使。}\par

% structure: p0018 source=h2 target=subsection reason=relative-heading-rank; base=h2

\subsection{第七章}

在塞农人地区有一个村庄，每年都遭受冰雹之灾。居民被极度的痛苦所迫，向玛尔定求助。由行政长官等级的奥司丕基乌斯（他的田地受风暴打击比别人的更严重）派遣了一个非常可敬的代表团到他那里。玛尔定在那里祈祷后，完全解除了整个地区流行的灾害，以致此后二十年他仍在肉身时，那些地方再无人遭受冰雹。但为免有人以为这事出于偶然，而非由玛尔定所致，当他去世的那一年，风暴重新降临那个地区。世界如此深切地感受到一位有信德之人的离世，正如它在他活着时合理地欢喜，也哀悼他的死亡。但若有听众信德薄弱，要求我们所说之事的见证人，我将提出的不是一个人，而是成千上万的人，甚至将召唤整个塞农人地区来为所经历的能力作证。但不谈这些，司铎雷弗里热里乌斯，我相信你记得，我们最近曾与我提到的那位奥司丕基乌斯的儿子罗穆禄斯（一位可敬而虔诚的人）谈论过此事。他把上述事情告诉我们，仿佛我们先前不知道似的；由于他害怕将来的收成持续受损，他正如你亲眼所见，带着许多哀叹惋惜玛尔定未能存活至今。\par

% structure: p0020 source=h2 target=subsection reason=relative-heading-rank; base=h2

\subsection{第八章}

但让我们回到\Person{阿维提亚努斯}：他在其他各处、在所有其他城市都表现出可怕的残酷，唯独在图尔未做任何伤害。是的，那以人血和不幸者的屠杀为养料的野兽，在真福者面前显得温和而和平。我记得\Person{玛尔定}有一天来到他那里，进入他的私室后，看见一个体型巨大的魔鬼坐在他背后。\Person{玛尔定}从远处向他吹气（如果我必须用一个几乎不是拉丁文的词来说的话），\Person{阿维提亚努斯}以为他在向\Emph{自己}吹气，喊道：\Quote{为什么，圣人，你这样对待我？}但\Person{玛尔定}说：\Quote{我不是向你，而是向那可怕地靠在你颈项上的那个。}魔鬼退却了，离开了它惯常的座位；众所周知，从那一天以后，\Person{阿维提亚努斯}变得温和了，或许是因为他现在明白自己一直是在遵行坐在他旁边的魔鬼的意志，或者是因为那被\Person{玛尔定}从座位上赶出的不洁之灵失去了攻击他的能力；而仆人羞于其主人，主人也不再强迫他的仆人。\par

在\Place{盎巴提恩塞斯}的一个村庄里，即一个古老的堡垒，如今大部分居住着弟兄们，你知道那里有一座用劳力建造的巨大偶像殿宇。那座建筑用最光滑的石头筑成，并配有塔楼；它呈圆锥形高高耸起，以其工程的庄严保存着那地方的迷信。真福者曾多次命令在那里定居为司铎的\Person{玛尔刻禄}摧毁它。过了一段时间返回时，他责备司铎，因为偶像殿宇的建筑物仍然矗立。司铎辩解说，如此巨大的结构即使由一队士兵或一大群公众的力量也难以推倒，更不用说\Person{玛尔定}认为靠软弱的圣职人员或无助的隐修士就能轻易做到。于是\Person{玛尔定}转而求助于他惯用的辅助力量，整夜守夜祈祷——结果，到了早晨，暴风骤起，将那偶像殿宇连根拔起。现在让这叙述依靠\Person{玛尔刻禄}的见证。\par

% structure: p0023 source=h2 target=subsection reason=relative-heading-rank; base=h2

\subsection{第九章}

在类似的工作中，我将利用另一个不逊于上述奇迹的事件，并有\Person{雷弗里盖里乌斯}的赞同。\Person{玛尔定}准备推倒一根巨大的柱子，柱顶立着一尊偶像，但无计可施。于是，他按照惯常的做法，转向祈祷。毫无疑问，那时一根与那柱子颇为相似的石柱从天而降，落在偶像上，将那似乎不可摧毁的整体碾为粉末：这还算小事，如果他只是以不可见的方式运用了天上的力量；但这些力量本身竟被人的眼睛看见，以可见的方式为\Person{玛尔定}服务。\par

再者，同一位\Person{雷弗里盖里乌斯}为我作证：一位患血漏的妇人，摸了\Person{玛尔定}的衣服，仿效福音中记载的那位妇人，立刻痊愈了。\par

一条蛇正穿河游向我们所在的岸边。\Person{玛尔定}说：\Quote{因主的名，我命令你回去。}圣者的话一出口，那毒物立即转身，在我们注视下游向对岸。我们都看出这并非没有奇迹；他深深叹息道：\Quote{蛇听我，但人不听。}\par

% structure: p0027 source=h2 target=subsection reason=relative-heading-rank; base=h2

\subsection{第十章}

\Person{玛尔定}习惯于在复活节期间吃鱼，在休息时间之前不久，他问是否准备好了。然后，负责隐修院外务且本人是熟练渔夫的执事\Person{卡托}告诉他，他整天一无所获，其他通常卖鱼的渔夫也同样一无所获。\Person{玛尔定}说：\Quote{去，放下你的钓线，必有收获。}正如\Person{苏尔皮基乌斯}在那里已描述过的，我们的住所靠近河边。于是，因为是节日，我们都去看我们的朋友钓鱼，大家都满怀希望，希望在\Person{玛尔定}本人的建议下为他获取鱼的努力不会落空。执事第一网就用一张很小的网捞出一条巨大的梭鱼，并欣喜地跑回隐修院，无疑带着某位诗人所表达的那种情感（因为我们用了一句有学问的韵文，既然我们是在与有学问的人交谈）——\Quote{并将他的俘虏野猪带给惊叹的\Place{阿尔戈斯}。}\par

确实，那位基督的门徒效法救主所行的奇迹——救主以这些奇迹作为榜样指示给他的圣人们——也显出了基督在他内工作，基督处处荣耀自己的圣人，将各种恩宠的恩赐赐予那一个人。御前侍卫\Person{阿尔波里乌斯}作证说，他看到\Person{玛尔定}在祭献时的手，仿佛装饰着最高贵的宝石，闪耀着紫色的光芒；当他移动右手时，他听到宝石相互碰撞的声响。\par

% structure: p0030 source=h2 target=subsection reason=relative-heading-rank; base=h2

\subsection{第十一章}

我现在要讲述一件事，他因时势而总是隐瞒，但他却无法对我们隐瞒。在这件事上，有一件神奇的事，就是一位天使曾与他面对面地交谈。皇帝\Person{玛克西穆斯}，虽然在其它方面无疑是个好人，但在\Person{普里西利安}被处死后，因某些司铎的建议而误入歧途。因此，他倚仗王权庇护了主教\Person{依塔基乌}——他是\Person{普里西利安}的控告者——以及他的其他同党，他们的名字无需提及。皇帝就这样阻止任何人控告\Person{依塔基乌}，因为藉着他的工具作用，一个什么样的人被判处死刑。现在，\Person{玛尔定}因许多遭受苦难者的严重事由被迫前往宫廷，他承受了在那里肆虐的风暴的全部力量。聚集在\Place{特里尔}的主教们被扣留在那座城市，他们每日与\Person{依塔基乌}交往，与他结成了同盟。当这个消息意外地传给他们说\Person{玛尔定}正在前来时，他们完全失去勇气，开始彼此咕哝和颤抖。恰好，在他们的影响下，皇帝已经决定派遣一些掌握绝对权力的护民官前往两西班牙省，搜捕异端者，找到后剥夺他们的生命或财产。毫无疑问，那场风暴也将摧毁大批真正的圣人，因为在不同类别的个人之间几乎不加区分。因为在这种情况下，判断仅凭外表作出，以至于一个人被认为异端，更多是因为他因恐惧而脸色苍白，或穿着某种特别的衣服，而不是因为他所宣认的信仰。主教们很清楚这样的行为绝不会令\Person{玛尔定}满意；但他们自知有罪，深感焦虑，生怕他到来后不与他们共融；他们深知，以他的榜样为引导，其他人也不会缺少，会跟随这样一位伟大人物的勇敢行动。因此，他们与皇帝策划了一个方案：派遣宫廷官员前去迎接\Person{玛尔定}，禁止他再靠近那座城市，除非他宣布与住在那里 的主教们保持和平。但他巧妙地挫败了他们的目的，宣布他将带着基督的平安到他们中间。最后，他在夜间进入，前往圣堂，只为祈祷。次日，他前往宫殿。除了他要呈递的许多其他请求（这些请求描述起来将十分冗长）之外，以下是主要的：为廷臣\Person{纳尔塞斯}和总督\Person{路卡狄乌}求情，两人都曾属于\Person{格拉提安}的派系，并且以超乎寻常的热忱——关于这一点现在不是详述的时候——因而招致了征服者的愤怒；但他的主要请求是，不要派遣掌握生杀大权的护民官前往西班牙。因为\Person{玛尔定}怀有虔诚的关怀，不仅要拯救那些地区的真基督徒（他们将在那次远征中受到迫害）脱离危险，甚至连异端者本人也要保护。但在第一天和第二天，狡猾的皇帝让圣人悬而不决，或者是为了向他强调此事的重要性，或者是因为他得罪了主教们而无法与他们和解，或者如同当时大多数人认为的，皇帝由于贪婪而反对他的愿望，贪图那些人的财产。因为我们听说，他确实是一个以许多杰出事迹著称的人，但在对抗贪婪方面他并不成功。然而，这可能是出于当时帝国的需要，因为国库已被前几任统治者耗尽；而他几乎经常期待内战或准备内战，因此他通过各种手段为帝国的防御寻求资源，这是很容易得到谅解的。\par

% structure: p0032 source=h2 target=subsection reason=relative-heading-rank; base=h2

\subsection{第十二章。}

与此同时，那些\Person{玛尔定}不愿与之共融的主教们惊恐地前往皇帝那里，抱怨说他们已被预先定罪；他们每人的\OriginalTerm{地位}{status}已经完了，如果\Person{玛尔定}的权威现在支持\Person{特奥格尼图斯}的固执——他是唯一一位至今公开宣判他们的人；这个人（指\Person{玛尔定}）不应该被接入城墙之内；他如今不仅是异端者的辩护者，更是他们的维护者；如果\Person{玛尔定}要扮演复仇者的角色，那么\Person{普里西利安}之死实际上什么也没有成就。最后，他们哭泣哀号，俯伏在地，恳求皇帝对这个人行使权力。皇帝几乎不得不也将\Person{玛尔定}判为异端者的下场。但尽管如此，虽然他倾向于过分偏袒主教们，但他并非不知道\Person{玛尔定}在信德、圣德和美德上超越所有凡人：因此他尝试另一种方法来战胜圣人。首先，他私下派人召来\Person{玛尔定}，以最和善的方式对他说，异端者是在公开审判的常规程序中被定罪的，而非由于司铎们的迫害；他没有理由认为与\Person{依塔基乌}及其同伙共融是一件该受谴责的事。他补充说，\Person{特奥格尼图斯}制造分裂，更多是由于个人仇恨，而非他所支持的事业；事实上，他是唯一一个暂时脱离共融的人；而其余的人并没有做出任何革新。他还提到，几天前召开的一次主教会议已裁定\Person{依塔基乌}不负有任何过失。当\Person{玛尔定}对这些话无动于衷时，皇帝便怒火中烧，急忙离开他的面前；与此同时，立即为\Person{玛尔定}为之求情的人们指定了刽子手。\par

% structure: p0034 source=h2 target=subsection reason=relative-heading-rank; base=h2

\subsection{第十三章。}

当此事传到\Person{玛尔定}耳中时，尽管已是夜晚，他仍赶往宫廷。他保证说：若这些人得以幸免，他就领受共融；只求将那些已被派往\Place{西班牙}摧毁教堂的护民官召回。毫不迟延：\Person{玛克西穆}允准了他所有的请求。次日，\Person{斐理克斯}被祝圣为主教的仪式正在安排中，此人无疑具有极大的圣德，确实堪当在更幸运的时代成为司铎。\Person{玛尔定}参加了那天的共融礼，认为暂时让步比不顾那些悬剑在颈之人的安全更为妥当。然而，主教们竭力劝说他签名以确认他共融的事实，但他无法被说服这样做。次日，他急忙离开那地方，在返回的路上，他充满悲哀和悔恨，后悔自己竟有一小时卷入那邪恶的共融。在离一个名为\Place{安代塔纳}的村庄不远处，那里幽深的树林绵延广阔，一片寂静，他坐下来，同伴们则稍往前走。他陷入深思，时而指责，时而辩护自己悲伤与行为的原因。忽然，一位天使站在他身旁说：\Quote{玛尔定，你感到内疚是公义的，但你无法以其他方式摆脱你的困境。重振你的德行，恢复你的勇气，免得你不仅现在危及你的名声，而且危及你自身的得救。}因此，从那时起，他小心提防，不再与\Person{依塔基乌}一派共融。但当他偶尔治愈附魔者时比平时更慢且恩宠较少时，他立刻含泪向我们承认，他感到自己的能力因那邪恶的共融而减弱，那共融他因急需且非出自乐意而短暂参与。此后他又活了十六年，但再也没有参加过任何主教会议，并小心避开所有主教集会。\par

% structure: p0036 source=h2 target=subsection reason=relative-heading-rank; base=h2

\subsection{第十四章}

但显然，正如我们所经历的，他以多倍的丰厚补偿了那暂时减弱的恩宠。后来我看到一个附魔者被带到他在隐修院门口；在那人触及门槛之前，他就被治愈了。\par

我最近听到一人作证说，当他航行在\Place{第勒尼安海}，沿着那条通往\Place{罗马}的航线时，暴风骤起，船上所有人都处于性命攸关的极度危险中。此时，一位尚未成为基督徒的埃及商人呼喊说：\Quote{玛尔定的天主，救救我们！}暴风雨立刻平息，他们得以继续所愿的航程，平静的大海继续保持完全的安宁。\par

\Person{里孔提乌}，一位信德之人，属于副将之列。当一场猛烈的疾病折磨着他的家人，病体躺满了他全屋，悲惨地证明了前所未闻的灾祸时，他通过一封书信恳求\Person{玛尔定}的帮助。此时，这位有福之人宣布所求之事难以获得，因为他灵里知道那家正在遭受天主的惩罚。然而，他没有停止连续七整天七整夜的祈祷和禁食，最终在祈求中达到了目标。\Person{里孔提乌}迅速体验到了天主的仁慈，飞跑到他那里，立即报告实情并感谢，说他全家已脱离一切危险。他还献出了一百磅银子，这位有福之人既不拒绝也不接受；但在这笔银子触及隐修院门槛之前，他已毫不犹豫地将其指定用于赎回俘虏。当弟兄们向他建议，应该留下一部分作为隐修院的费用，因为他们全体都难以获得必要的食物，且许多人急需衣服时，他回答说：\Quote{教会既喂养又穿戴我们，只要我们不以任何方式显得是自己供给了自己的需要。}\par

此刻，我想起了那位杰出之人的许多奇迹，这些奇迹我们更容易赞叹而非叙述。你们无疑都承认我所说的是事实：他的许多事迹无法用言语表达。例如，下面这一件事，我认为我们无法完全按原样叙述。有一位弟兄（你们并非不知道他的名字，但必须隐藏他的身份，以免使一位虔诚的人蒙羞）——我说，有一位弟兄，发现炉子有大量煤炭，便拉过一张凳子，在炉火旁叉开腿坐着，暴露了身体。\Person{玛尔定}立刻察觉到在神圣的屋顶下发生了不当之事，大声喊道：\Quote{谁暴露身体，玷污了我们的居所？}那弟兄听到这话，凭良心意识到自己正是被责备的人，立刻几乎昏厥地跑到我们这里，承认了自己的羞耻；然而，这完全是藉着\Person{玛尔定}权能的彰显而成就的。\par

% structure: p0041 source=h2 target=subsection reason=relative-heading-rank; base=h2

\subsection{第十五章}

\Quote{有一次，他坐在他的那个木制座位上（你们都知道的），放在围绕他住所的小露天庭院里，他看见两个魔鬼坐在高悬于隐修院上方的岩石上。然后他听见他们用急切而快活的语调发出如下邀请：\InnerQuote{过来，\Person{布里克提奥}，过来，\Person{布里克提奥}。}我相信他们远远地看见那个可怜的人走来，意识到他们在他心中激起了多么大的精神狂乱。没有耽搁：\Person{布里克提奥}怒气冲冲地冲了进来，在那里，充满疯狂，他对\Person{玛尔定}倾泻出千种指责。因为前一天他受到了\Person{玛尔定}的责备，因为他从前在进入圣职之前一无所有，实际上是由\Person{玛尔定}自己在隐修院中抚养长大的，现在却养着马匹并购买奴隶。因为当时，许多人控告他不仅购买了蛮族的男童，而且还购买了容貌秀丽的少女。这个可怜的人因这些事而痛苦愤怒，并且我相信主要是受那些魔鬼的冲动所驱使，向\Person{玛尔定}发起如此猛烈的攻击，以至于几乎克制不住要动手打他。圣人方面，带着平静的面容和安宁的心，试图用温和的话语来制止这不幸之人的疯狂。但邪恶的灵在他内如此得胜，以至于连他自己本已十分虚妄的理智都不受控制。他嘴唇颤抖，面色多变，因愤怒而苍白，吐出罪恶的话语，声称自己比养育他的\Person{玛尔定}更圣洁，因为他从幼年起就在隐修院中，在教会的神圣体制中长大，而\Person{玛尔定}最初，正如他不能否认的，曾受到军人生活的玷污，现在又因他那毫无根据的迷信和对异象的荒谬幻想而完全陷入昏聩。在他说了许多这样的话，以及其他更加尖刻的话（这些最好不提）之后，当他的愤怒得到满足时，他离开时似乎感到自己已经完全为自己的行为辩护了。但他沿着出去的路快步冲回，我相信此时魔鬼已因\Person{玛尔定}的祈祷而从他心中被逐出，他现在被带回悔改。他很快回来，扑倒在\Person{玛尔定}脚前，请求宽恕并承认自己的错误，同时，终于恢复了更好的心智，他承认自己曾受到魔鬼的影响。\Person{玛尔定}宽恕恳求者并不困难。然后圣人向他并向我们所有人解释，他如何看见他被魔鬼驱赶，并声称他并不因那些堆到他身上的指责而动摇；因为这些指责实际上反而伤害了说话的人。后来，当同一\Person{布里克提奥}多次在他面前被控告犯下许多重罪时，\Person{玛尔定}无法被说服免除他的司铎职，以免他被怀疑是报复对自己的伤害，而他经常重复这句话：\InnerQuote{如果基督容忍了\Person{犹达斯}，我为什么不能容忍\Person{布里克提奥}呢？}}\par

% structure: p0043 source=h2 target=subsection reason=relative-heading-rank; base=h2

\subsection{第十六章}

\Quote{这时，\Person{波斯图米亚努斯}喊道：\InnerQuote{让我们附近的那位知名人士听听这个榜样，他明智时，既不注意现在也不注意将来，但如果被冒犯，就会陷入完全愤怒，无法自控。然后他对神职人员发怒，猛烈攻击平信徒，并为自己的报复搅动整个世界。他会在这种争辩中持续三年不停息，拒绝因时间或理性而软化。这个人的状况值得哀悼和同情，即使这只是他受困扰的无可救药的恶疾之一。但是，我的高卢朋友，你应当经常向他心中提醒这类忍耐与平安的榜样，好让他知道如何发怒，也知道如何宽恕。如果他不巧听到我这番话——它被简短地插入我们的谈话中，且是针对他自己——让他知道，我说这些话，与其说是以敌人的口吻，不如说是以朋友的心意；因为我希望，如果可能的话，人们说他像\Person{玛尔定}主教，而不是像暴君\Person{法拉里斯}。但我们还是离开他吧，因为提到他并不令人愉快，让我们回到我们的朋友\Person{玛尔定}那里，高卢人啊。}}\par

% structure: p0045 source=h2 target=subsection reason=relative-heading-rank; base=h2

\subsection{第十七章}

那时我见夕阳西下，黄昏已近，便说：\Quote{日子已尽，\Person{波斯图米亚努斯}，我们该起身了；同时，也该给这些如此热心的听众一些休憩。至于\Person{玛尔定}，你不该指望谈论他会有什么尽头：他太过广博，无法在任何谈话中被完全包容。此刻，你将把你所听闻的关于这位杰出人物的事迹带到东方；当你重访昔日旧地，途经各海岸、各处、各港口、各岛屿、各海洋时，务必将\Person{玛尔定}的名声与荣耀传扬给万民。尤其要记住，不可遗漏\Place{坎帕尼亚}；即使你的路线偏离常道，任何耽搁造成的耗费，也不及你前往那一区拜访\Person{帕乌利努斯}——这位举世闻名、备受赞誉的人物——来得重要。我恳请你首先向他展露我们昨日完成或今日所论的整卷谈话。你将向他陈述一切，向他复述一切；以便届时通过他，\Place{罗马}得以了解此人的神圣功德，正如他之前将那本我们的小书不仅传遍\Place{意大利}，甚至传遍整个\Place{伊利里亚}。他并不嫉妒\Person{玛尔定}的荣耀，而是以最虔诚的心钦佩\Person{玛尔定}在\Person{基督}内圣德的高超，他不会拒绝将我们的领袖与他自己的朋友\Person{斐理克斯}相比。其次，你若偶然渡海到\Place{阿非利加}，便将所闻禀告\Place{迦太基}；虽然，如你所说，它已认识此人，但如今更应深入了解他，以免它只仰慕自己的殉道者\Person{西彼廉}，虽则他已用自己的圣血祝圣了自己。而后，若你稍向左拐，进入\Place{阿哈雅}湾，务使\Place{科林托斯}和\Place{雅典}知道：学园中的\Person{柏拉图}并不比\Person{玛尔定}更智慧，监狱中的\Person{苏格拉底}并不比\Person{玛尔定}更勇敢。你将对他们说：\Place{希腊}确实有福，堪当聆听一位宗徒的辩护，但\Person{基督}绝未舍弃\Place{高卢}，因为他恩赐它拥有\Person{玛尔定}这样的人。当你远至\Place{埃及}，虽然它因自己众多圣人的数量和德行而正当地自豪，但仍不应不屑于听闻：\Place{欧洲}如何不逊于它或整个\Place{亚洲}，只因它拥有\Person{玛尔定}一人。}\par

% structure: p0047 source=h2 target=subsection reason=relative-heading-rank; base=h2

\subsection{第十八章}

\Quote{但你若再从那里启航前往\Place{耶路撒冷}，我嘱托你一件关乎我们忧伤的事：你若来到著名的\Place{托勒迈}海岸，务必仔细打听我们的那位朋友\Person{彭波尼乌斯}葬于何处，不要拒绝在那异乡的土地上凭吊他的遗骸。在那里洒下许多眼泪，既出于你自己的情感，也出于我温情的眷爱；虽然只是微薄之礼，仍请用紫花和香草撒遍那土壤。你要对他说——不是粗暴，不是严苛，而是带着同情的话语，而非责难的口气——倘若他当年肯听从你一时，或始终听从我，并且邀请\Person{玛尔定}而非我不愿提名的那人，他绝不会这样残忍地与我分离，也不会被一堆陌生的尘土掩埋，在海上遭遇海盗般的厄运，死于船难，勉强在遥远的海岸得以下葬。愿那些为了报复他而图谋伤害我的人，将这一切视为自己的作为；愿他们目睹自己的荣耀，既已报仇，从此不再攻击我。}\par

我们以极其哀伤的语气说出这些悲恸之词，众人的泪水因我们的哀叹而涌出，终于辞别而去，心中固然对\Person{玛尔定}满怀深深的钦仰，却也不亚于因自己的悲伤而生的哀痛。\par

\SourceRecord{Translated by Alexander Roberts.；Nicene and Post-Nicene Fathers, Second Series；Vol. 11.；Edited by Philip Schaff and Henry Wace.；Buffalo, NY: Christian Literature Publishing Co.,；1894.；<http://www.newadvent.org/fathers/35033.htm>.}
